% ============================================================
%  Chapter 03: Sociology as Science: Research Methodology,
%  Positivism, Non-Positivism & Critique of Scientific Method
% ============================================================

\chapter{Sociology as Science: Methodologies, Positivism \& Critiques}

\section{Answer Writing Practice and the Foundations of Reasoning}

\subsection{Setting the Frame: Science as a Way of Life}
Unit 1 is now complete . Unit 2 is mainly about how you think---it is about you . The unit opens with a single question: what do we understand by science? Is science only about those laboratory experiments we have been seeing since childhood? 
\begin{itemize}
    \item Laboratory work---the white rat, laser radiation, and similar images---is scientific experiment . It is not the whole of science .
    \item People generally say religion is a way of life . Science is also a way of life . The question that follows is: how does science govern our lives, every single day? 
\end{itemize}

\begin{KeyConcept}
\begin{itemize}
    \item Science is not confined to the laboratory . Like religion, it is a way of life---a way of thinking that we practise daily without noticing it .
    \item The laboratory image (white rat, laser radiation) belongs to scientific experiment, which is only one part of science .
\end{itemize}
\end{KeyConcept}

\begin{QuickRevision}
\begin{itemize}
    \item Unit 2 is about how you think .
    \item Science is a way of life, not only laboratory experiment .
    \item Central question of the unit: how does science govern our everyday lives? 
\end{itemize}
\end{QuickRevision}

\subsection{The Classroom Exercise: ``Is Sociology Common Sense?''}
Before entering Unit 2, an answer-writing exercise was set in class . Students were asked to take out one A4 sheet and write the question: \textit{``Is sociology common sense?''} 
\begin{itemize}
    \item The instruction was not to think about the answer as it will be after Unit 2 . Students had to write with the existing, preliminary knowledge available from Unit 1 alone .
    \item Once Unit 2 is done, a full-fledged understanding will develop . The purpose here was to check how the question would be solved with the knowledge already in hand .
    \item It is a 10-mark question, which means the answer must not exceed 150 words . The exercise was timed, and roughly 8 minutes were available .
\end{itemize}

\begin{PYQLink}
\begin{itemize}
    \item \textit{``Is sociology common sense?''}---treated as a 10-mark question (150-word limit) .
    \item The same question will be returned to after Unit 2 and again after Unit 4, to show how the content of the answer changes as knowledge grows .
\end{itemize}
\end{PYQLink}

\begin{QuickRevision}
\begin{itemize}
    \item Exercise question: is sociology common sense? 
    \item 10 marks $\rightarrow$ 150 words, written in about 8 minutes .
    \item Answer to be attempted with Unit 1 knowledge only .
\end{itemize}
\end{QuickRevision}

\subsection{Step 1 and 2 --- Identify the Keywords, Then Define Them}
The approach to the question was demonstrated step by step . The order matters: you do not read the question first .
\begin{enumerate}
    \item \textbf{Identify (encircle) the keywords:}  Here the keywords are \textit{common sense}, \textit{sociology}---and also \textit{is} . In effect the whole question is a keyword . ``Is A equal to B?'' is again a question of causation .
    \item \textbf{Define the keywords:} 
    \begin{itemize}
        \item \textit{Sociology:} Remember, there is no one definition of sociology . The line written was: sociology is a scientific study of society and human behaviour . Since it is a 10-mark question, the definition is kept to very few lines .
        \item \textit{Common sense:} A set of knowledge characterised by opinions, assumptions, biases or values .
        \item Unit 2, which is all about what makes something scientific, has not yet been covered . Even so, with a basic understanding of science one can say that common sense is not scientific---because opinions, assumptions and biases are not scientific, even if the exact meaning of ``scientific'' is not yet known .
    \end{itemize}
    \item \textbf{Now read the question:}  Having defined both A and B, the question asks you to draw a relation between them; only after drawing that relation can you say whether they are equal or not .
\end{enumerate}

\begin{CriticalPoint}
\begin{itemize}
    \item The sequence is keywords $\rightarrow$ definitions $\rightarrow$ read the question, not the other way round .
    \item ``Is A equal to B?'' is a question of causation; both terms must be defined before a relation can be drawn .
    \item Honest use of limited knowledge is acceptable: you may say common sense is not scientific even while admitting the exact meaning of scientific is yet to be studied .
\end{itemize}
\end{CriticalPoint}

\begin{QuickRevision}
\begin{itemize}
    \item Three steps: identify keywords, define keywords, then read the question .
    \item Sociology = scientific study of society and human behaviour (no single definition exists) .
    \item Common sense = set of knowledge characterised by opinions, assumptions, biases or values---not scientific .
\end{itemize}
\end{QuickRevision}

\subsection{Building the Body: Sociology as a Scientific Study}
With sociology defined as scientific and common sense as not scientific, the first conclusion that arises is that both are opposite . The examiner has not stated that sociology is common sense; he has only asked . So the answer moves forward line by line .

\paragraph{The reasoning chain written in the answer}
\begin{enumerate}
    \item Sociology, at least in its initial phase, aimed to study society as a system governed by laws (third line, after the two definitions) .
    \item Auguste Comte developed the law of three stages . Early sociologists like Comte held that the social world, just like the natural world, is governed by laws .
    \item He was heavily inspired by Enlightenment scholars / the Enlightenment era, which believed in the study of society with reason and experimentation .
\end{enumerate}
Everything is thus in sync---each sentence follows from the one before it . The answer is not haphazard .

\paragraph{André Béteille and the new scholars}
\begin{itemize}
    \item Study of society with reason and experimentation is the core of sociology . From there the answer moves to newer scholars .
    \item \textbf{André Béteille:} Sociology is a systematic study governed by data collection, classification, cause--effect and comparison .
    \item He also believed that sociologists study the function of an institution to the larger system .
\end{itemize}

\begin{CriticalPoint}
\begin{itemize}
    \item At this point the word limit begins to press . Data collection, classification, cause--effect, reason, experimentation, observation cannot all be expanded---some importance must still be given to common sense .
    \item The deliberate decision taken in class: stop here on ``sociology as scientific study'' and move on .
\end{itemize}
\end{CriticalPoint}

\paragraph{Why these features separate sociology from common sense}
\begin{itemize}
    \item An individual does not have the ability to compare seemingly different institutions; and even if he compares, he cannot arrive at a generalisation that a pattern runs through them .
    \begin{itemize}
        \item The comparative method is not something an ordinary individual develops . He has not developed the critical thinking that allows him to compare school with country, or caste with class . To him these are incomparable entities---``you cannot compare caste and class, both are different''; ``you cannot compare religion, race and caste'' .
    \end{itemize}
    \item A common-sensical person cannot see that each institution, just like each organ in the human body, has a function in the maintenance of the whole .
    \begin{itemize}
        \item Each organ has a function in maintaining the body; similarly each institution contributes to the overall functioning of society . This is not how a common-sense man or woman thinks .
    \end{itemize}
    \item One line then closes this half of the answer: \textit{``These are not found in common sense thinking.''}  Sociological thinking is very different from common sense thinking .
\end{itemize}

\begin{ExampleBox}
Organic analogy used informally: each organ of the human body has a function in maintaining the body; each institution similarly contributes to the functioning of society---a link a common-sense observer does not make .
\end{ExampleBox}

\begin{CriticalPoint}
\begin{itemize}
    \item Underline important keywords as you write, to make the answer visible to the examiner .
    \item Do not use black, pink or red pens . A blue pen is more than sufficient---simply underline what you have already written with it .
\end{itemize}
\end{CriticalPoint}

\begin{QuickRevision}
\begin{itemize}
    \item Initial-phase sociology: society as a system governed by laws---Comte, law of three stages, Enlightenment inspiration .
    \item André Béteille: sociology = systematic study via data collection, classification, cause--effect, comparison; institutions studied for their function to the larger system .
    \item Common sense lacks the comparative method, generalisation, and the sense of institutional function .
    \item Answer discipline: underline keywords with the blue pen only .
\end{itemize}
\end{QuickRevision}

\subsection{The Turn: Writing ``However'' and Bringing in Social Action}
So far the answer has established that sociology is not at all like common sense---an indirect ``No'' . But the answer cannot stop at a yes or a no . Both have to be written .

\begin{KeyConcept}
\begin{itemize}
    \item When you have to shift to another line of thinking---that is, to reject what you have just argued---you must begin with one word: \textbf{``However''} .
    \item The moment you write ``however'', you force yourself to write something opposite to what you have written . You will automatically feel the force, because you cannot repeat the same thing after it .
\end{itemize}
\end{KeyConcept}

\paragraph{The ``however'' half of the answer}
\begin{itemize}
    \item However, later sociologists argued that human behaviour is highly complex and cannot be theorised . You cannot make a theory out of human behaviour; you cannot predict what a society or a person is going to do next .
    \item And therefore social action should be studied without bothering about laws or theories---that is, why an individual does whatever he does .
    \item The one scholar available for social action at this stage is \textbf{Max Weber} . With Unit 1 knowledge we know only that social action means the study of meanings and motives of individuals, given by Weber---and nothing else .
    \item Concluding line of this half: the study of social action brings sociology closer to common sense .
\end{itemize}

\begin{ExampleBox}
\textbf{The burning building:} A building catches fire . Some people run away . Others, instead of running away, jump to the third floor and try to rescue those inside, risking their own lives . You may theorise that people will run away---but then how do you explain those who climb up without thinking about their own lives? Human behaviour cannot be scientifically theorised in this way .
\end{ExampleBox}

\begin{CriticalPoint}
\begin{itemize}
    \item Initial sociologists were hell-bent on developing laws . Later, Weber realised that theorising a law out of human behaviour was misplaced, and focused on the individual---on the meanings and motives we attribute to our actions and to institutions .
    \item It is not only about institutions contributing to the functioning of society; individuals also look at institutions in a very unique fashion, which the earlier writers did not study .
\end{itemize}
\end{CriticalPoint}

\begin{QuickRevision}
\begin{itemize}
    \item ``However'' is the pivot word that forces the counter-argument .
    \item Later sociologists: human behaviour is highly complex and cannot be theorised .
    \item Hence social action---Weber---study of meanings and motives, without bothering about laws .
    \item Social action brings sociology closer to common sense .
\end{itemize}
\end{QuickRevision}

\subsection{Writing the Conclusion and the Lesson on Form over Content}
Having written both ``no'' and ``yes'', the answer cannot leave the examiner without a position . A both-yes-and-no answer reads as an inability to decide---and, as was put in class, an administrator who cannot make decisions will not be selected . A balanced conclusion is therefore required .

\begin{KeyConcept}
\begin{itemize}
    \item \textbf{Conclusion line as dictated:} \textit{In conclusion, the very journey of the scope of sociology, ranging from the study of institutions to the study of social action, reflects the growth of the discipline just like a human body.} 
    \item Our body matures after a certain point of time; sociology also matured later on . Initially it insisted on one thing only---the childhood insistence on one particular chocolate---and later it grew out of that insistence .
    \item Just as human behaviour matures, a discipline also matures . That is why there is ever-evolving research in every discipline, and every year people who have contributed immensely to the growth of a discipline are rewarded---for instance, those who receive the Nobel Prize in one area or another .
    \item The whole trajectory shows nothing but a journey of maturity . In this way you have evaded a flat yes or no and still reflected a mature position .
\end{itemize}
\end{KeyConcept}

\paragraph{What was deliberately left out --- and why}
\begin{itemize}
    \item The answer did not talk about Zygmunt Bauman, defamiliarisation, or several other scholars who could have been introduced .
    \item When a countdown is running and so much information has to be compressed into 8 minutes, your brain automatically starts filtering what is important and what is not .
    \item This filtering depends on revision :
    \begin{itemize}
        \item \textit{Revised once:} You cannot develop a structure and may start writing haphazardly .
        \item \textit{Revised multiple times:} You feel confusion: should this scholar go here, is it relevant, or should I skip it? 
        \item \textit{Third stage:} The mind is absolutely clear about what the structure should be; whatever comes to mind, you write . At this stage you do not think about content, you think about the form of the content .
    \end{itemize}
\end{itemize}

\begin{KeyConcept}
\begin{itemize}
    \item Just as formalist scholars talk more about the form of a relationship than its content, an answer is judged on form .
    \item If the form---presentation and clarity of thought---is reflected in what you have written, you will get marks irrespective of which scholar you used, unless the scholar used is completely wrong .
    \item \textbf{Structure is more important than content} .
\end{itemize}
\end{KeyConcept}

\begin{CriticalPoint}
\begin{itemize}
    \item This exercise makes you realise on your own that there is no point in reading multiple sources; it only creates more difficulty .
    \item \textbf{The formula stated:} $\text{Limited information} + \text{Unlimited revision} = \text{Clarity of thought}$ .
    \item The same question will be revisited after Unit 2, and again after Unit 4 . Each time, the content of the answer changes---showing that the more you learn, the more your thinking and your approach to the question change .
    \item \textbf{On working in teams:} Although you are competitive with each other, at this stage you should cooperate . Competition can come after clearing the examination; without friends pointing out your mistakes, improvement is limited, since a teacher cannot always be present .
\end{itemize}
\end{CriticalPoint}

\begin{QuickRevision}
\begin{itemize}
    \item Never end on a bare yes or no---give a balanced conclusion .
    \item Conclusion device used: the journey of sociology from institutions to social action = maturity of a discipline, like a maturing human body .
    \item Structure/form $>$ content in answer writing .
    \item Limited information, unlimited revision = clarity of thought .
\end{itemize}
\end{QuickRevision}

\subsection{Entering Unit 2 --- Syllabus Map and Scope of the Discipline}
The name of the unit is \textbf{Sociology as Science} . The first topic in the syllabus is science, scientific method and critique---so it is necessary first to understand what science is, then the scientific method, and only then to criticise the scientific method .
\begin{itemize}
    \item In Unit 1 there was a great deal of talk about observation, reasoning, experimentation . Scientific method is a term already encountered while reading the handouts, and so is empiricism---but only as a preliminary idea . All these words will now be tackled .
    \item Social action theory, structural functionalism and Marxism also belong to this unit; a preliminary idea of each already exists .
\end{itemize}

\paragraph{Sociology versus other disciplines --- what will and will not be covered}
\begin{itemize}
    \item For anthropology versus sociology, read NCERT Class 11, \textit{Introducing Sociology}, Chapter 1---the comparison is contained in about one page .
    \item Philosophy and political science comparisons will not be taken up separately . In earlier batches they were explained, but the input--output ratio is poor: hardly a 10 or 20 mark question comes from this area, once in a blue moon .
    \begin{itemize}
        \item \textit{Sociology versus philosophy} has never been asked by UPSC . Philosophy is the mother of all sciences---a position attributed in class to the scholar \textbf{M.~T. Cicero}---and it would be unfair to ask you to compare the son with the mother .
        \item \textit{Political science:} Such questions used to be asked long ago . Today the interest is in how you understand the larger discipline and how sociology and other disciplines interact with each other .
    \end{itemize}
    \item \textbf{Interdisciplinarity} is more important today, and will be dealt with in Unit 2 . A student of a humanities university has to take one more discipline besides the main one---a sociology postgraduate may have to study environmental science or psychology in one semester or another . We no longer talk about conflict between disciplines .
\end{itemize}

\begin{PYQLink}
\begin{itemize}
    \item Sociology versus philosophy---never asked by UPSC .
    \item Sociology versus political science---asked long ago; today the framing has shifted to interdisciplinarity and disciplinary interaction .
\end{itemize}
\end{PYQLink}

\begin{QuickRevision}
\begin{itemize}
    \item Unit name: Sociology as Science; first sub-topic: science, scientific method and critique .
    \item Unit 2 also carries social action theory, structural functionalism and Marxism .
    \item Anthropology vs sociology $\rightarrow$ NCERT Class 11 \textit{Introducing Sociology}, Chapter 1 .
    \item Philosophy/political science comparisons skipped; interdisciplinarity is the current framing .
\end{itemize}
\end{QuickRevision}

\subsection{Reasoning in Everyday Life: The Car That Would Not Start}
Before science can be defined, reasoning has to be understood . Reasoning is not something found in a laboratory; it is something all of us do every day, at home and outside . Student responses in the chat box described it as the cause of something, ``everything has a cause'', rational thought, employing rational thinking, an analytical approach . In earlier times, the cause of the way we organised our lives was religion .

\begin{ExampleBox}
\textbf{The car that will not start:}
\begin{itemize}
    \item I have a car; I try to switch it on and it does not turn on . The moment I observe this, I begin to reason: why? 
    \item The explanations---the so-called reasons---that come to mind: (1) the battery is dead; (2) the car has run out of fuel; (3) the engine is dead; (4) I am sitting in the wrong car---it looks so much like mine; (5) the keys in my hand are not the right ones .
    \item This whole process of trying to work out why the car is not functioning is nothing but reasoning .
    \item The first thought was that the battery must be dead . You do not consider everything at once: you start with one explanation, treat it as true, and work on it .
    \item So the second step is: change the battery .
    \item Third question---why did I change the battery? Because I believed the problem lay with the battery . In doing so I am also making a prediction: if I change the battery, everything will be right . The whole process implies prediction .
    \item Whether that reason is right or wrong is yet to be tested . You have assumed it to be the cause---and that assumption is nothing but a hypothesis .
\end{itemize}
\end{ExampleBox}

\begin{KeyConcept}
\begin{itemize}
    \item \textbf{Hypothesis (as dictated):} Any statement or any assumption which is yet to be tested .
    \item Five possible reasons behind the car not working = five hypotheses (hypothesis 1 to hypothesis 5) .
    \item The battery is changed---and the car still does not function . That means the hypothesis you went with was wrong; the cause is something else . When the assumption was tested, it turned out not to be a valid reason .
    \item You then move to the next reason: petrol . This is the \textbf{alternative hypothesis}---the alternative reason . The car is pushed to the petrol station with the help of friends, the tank is filled full---and the car still does not function .
\end{itemize}
\end{KeyConcept}

\begin{QuickRevision}
\begin{itemize}
    \item Reasoning = working out the possible causes of an observed problem .
    \item Car example: five possible causes = five hypotheses .
    \item Hypothesis = any statement or assumption yet to be tested .
    \item A failed test $\rightarrow$ move to the alternative hypothesis .
\end{itemize}
\end{QuickRevision}

\subsection{Hypothesis, Prediction, Experimentation --- The Steps of Reasoning}
The same structure was then demonstrated with a problem the students themselves had faced, to show that this is how we think in everyday life . Identification of the problem comes first; only then does reasoning begin .

\begin{ExampleBox}
\textbf{The live stream problem:}
\begin{itemize}
    \item The problem identified: the live stream was not working properly .
    \item Possible causes, each of them a hypothesis because none is yet tested: (1) something is wrong at the technical team's end; (2) my internet connection could be slow; (3) something is wrong with my device; (4) something is internally wrong with the internet service provider in that region---isolated places such as the North-East or Jammu \& Kashmir, where users feel lag; (5) the browser is not the right one .
    \item If the ISP hypothesis is taken as the cause, the corrective step is to change the ISP or buy higher bandwidth---which itself implies the prediction that the issue will then stop . The change is made; the problem persists; the hypothesis is nullified .
    \item Because the hypothesis has been nullified, the technical team would not agree that the problem is at their end . A hypothesis that has been nullified is called the \textbf{null hypothesis} .
    \item When you nullify a hypothesis, the next step is to think of the alternative hypothesis .
\end{itemize}
\end{ExampleBox}

\paragraph{The dictated definition of reasoning}
The concept of reasoning consists of :
\begin{enumerate}
    \item Development of hypothesis 
    \item Prediction 
    \item Experimentation 
    \item Confirmation or rejection of the hypothesis 
\end{enumerate}
Changing the ISP is an experiment . Changing the battery is an experiment . Both are done to see whether the prediction is correct or not---so once you predict, you also experiment .

\begin{CriticalPoint}
\begin{itemize}
    \item Do we always think like this? No . Sometimes yes, many times no .
    \item When the car is not functioning you follow these steps . When you are not functioning---a poor performance in a test, a sports event, a music competition---you do not .
    \item Instead of introspection, you start listening to toppers, relying on what others have to say . When the car failed you did not search for a video; you called a mechanic or did it yourself .
    \item To doubt is part of human nature---but the doubt raised by poor performance is usually answered by others' opinions rather than by testing .
\end{itemize}
\end{CriticalPoint}

\begin{QuickRevision}
\begin{itemize}
    \item Steps of reasoning: hypothesis $\rightarrow$ prediction $\rightarrow$ experimentation $\rightarrow$ confirmation or rejection .
    \item A hypothesis that fails the test = null hypothesis; the next one taken up = alternative hypothesis .
    \item Problem identification always precedes reasoning .
\end{itemize}
\end{QuickRevision}

\subsection{Reasoning in History: The Middle Ages and the Break From It}
The question was then widened from the individual to humanity . Did humanity follow these steps in the middle ages? Nobody used to experiment to check whether what the church was telling them was right or wrong .
\begin{itemize}
    \item The church held that life on the earth is inconsequential; one should be ready to go to heaven; for that one must come to church every day, pay heed to what is taught, and be the child of God . This temporary existence given by God is something to be thankful for .
    \item Suppose I am told I am destined to go to heaven . Can I test it? If I had no doubt and did not question, then there is no point in talking about hypothesis---and with no hypothesis there is no prediction and no experimentation, which are the later stages .
\end{itemize}

\begin{CriticalPoint}
\begin{itemize}
    \item The chain fails at the first link: no doubt $\rightarrow$ no hypothesis $\rightarrow$ no prediction $\rightarrow$ no experimentation .
    \item A claim such as ``you are destined to go to heaven'' cannot be put to any test at all .
    \item It is only during the middle ages, or right after them, that some charismatic personalities dared to challenge this notion and talked about testing it---about experimentation, prediction and the rest .
    \item These are the things that make us what we are: the man who doubts, who reasons---a \textbf{rational being} . Someone who thinks, doubts, tests the given information, experiments, predicts and observes, all by himself .
\end{itemize}
\end{CriticalPoint}

\begin{ExampleBox}
\begin{itemize}
    \item Names cited: Copernicus, Galileo Galilei, Isaac Newton---that is why they became famous .
    \item And likewise Hobbes, Locke, Rousseau, Voltaire, Mary Wollstonecraft, Adam Smith .
    \item They too had multiple hypotheses, predicted and experimented with them, and finally said, for instance, that capitalism is better than mercantilism, or that the sun is at the centre, not the earth---observed through the telescope, with multiple calculations to support it .
    \item The point: they used to predict and observe everything by themselves .
\end{itemize}
\end{ExampleBox}

\begin{KeyConcept}
\begin{itemize}
    \item \textbf{Dictated line:} \textit{``The essence of science lies in reasoning.''} 
    \item It follows that to be scientific means to have a hypothesis that can be tested, that a prediction can be made on, and that can be experimented on .
\end{itemize}
\end{KeyConcept}

\begin{QuickRevision}
\begin{itemize}
    \item Middle ages: religious explanation, no doubt, no testing---hence no hypothesis, prediction or experimentation .
    \item The break came with figures who insisted on testing: Copernicus, Galileo, Newton; Hobbes, Locke, Rousseau, Voltaire, Wollstonecraft, Adam Smith .
    \item The essence of science lies in reasoning .
\end{itemize}
\end{QuickRevision}

\subsection{Scientific Method and the Test of Testability}
What was covered in this session was described as an itinerary for Unit 2---the basic ingredients you must have before the actual cooking begins .

\begin{ExampleBox}
\textbf{Itinerary analogy:} When you travel to a place such as Mussoorie, you know what luggage should be with you and which spots you will visit---the whole travel plan is your itinerary . Today's class is the itinerary for Unit 2; the basic ingredients before you make the dish .
\end{ExampleBox}

\begin{KeyConcept}
\textbf{Scientific method (as defined in class):} A method which involves systematic steps to carry out experiment and to reach the truth .
\begin{itemize}
    \item The car example was exactly this: a systematic plan, prediction, and each step cross-checked and experimented on . Many predictions failed, and finally, when all the steps had been nullified, a new car had to be bought .
    \item That was scientific method carried out without the white coat a scientist wears in a lab . One can observe, predict, experiment and finally reach the truth without being a scientist . So do not be afraid of the words science or scientific method---it is simply a way of life .
    \item Unfortunately this way of life was absent in the middle ages, when religion was the only way of life and intellectual thinking was highly religious; everything was explained in a religious fashion .
\end{itemize}
\end{KeyConcept}

\begin{ExampleBox}
If a car had failed to start in the middle ages, the explanation would have been: it is the curse of God . This too is a hypothesis, since it is yet to be tested . The problem is that you can never test whether you have been cursed---whereas battery, engine, fuel and the car itself can all be tested .
\end{ExampleBox}

\begin{KeyConcept}
\begin{itemize}
    \item \textbf{Dictated criterion:} \textit{In order for a statement, a theory or a hypothesis to be called scientific, it should be testable.} 
    \item Anything we speak in everyday life can be scientific---only and only if it is testable .
\end{itemize}
\end{KeyConcept}

\begin{ExampleBox}
\textit{``I am coughing a lot; I think I have contracted coronavirus.''}  Is this scientific? Yes---because it can be tested . That was one hypothesis . A second hypothesis was also in mind: it may be casual---a cold drink taken over something hot, a natural reaction of the body, nothing to worry about . The test at the clinic found no coronavirus . But that does not by itself mean you do not have it---the problem may lie with the experimentation, which is a further layer, deliberately left for later .
\end{ExampleBox}

\begin{KeyConcept}
A theory, statement or hypothesis can be tested through two methods :
\begin{enumerate}
    \item \textbf{Induction} or inductive reasoning 
    \item \textbf{Deduction} or deductive reasoning 
\end{enumerate}
\end{KeyConcept}

\begin{QuickRevision}
\begin{itemize}
    \item Scientific method = systematic steps to carry out experiment and reach the truth .
    \item You do not have to be a scientist, or be in a laboratory, to follow it .
    \item Testability is the criterion: a statement, theory or hypothesis is scientific only if it can be tested .
    \item Two methods of testing: induction and deduction .
\end{itemize}
\end{QuickRevision}

\subsection{Inductive Reasoning and the Problem of Hasty Generalisation}

\begin{ExampleBox}
\textbf{The white swan story:}
\begin{itemize}
    \item Suppose I am a European . One fine day I go out of my house, sit by a lake and see beautiful swans---and when the word swan is spoken, a white swan is what comes to mind .
    \item I photograph them, and over the next days, months and up to the next year I keep seeing white swans---in England, in Norway, everywhere . After about 100 swans, I develop the generalisation: \textit{all swans are white} .
    \item I have gone from the specific (100 swans actually seen) to the general (all swans) .
    \item Then one day I go to New Zealand---or Australia---and see a black bird which I am told is also a swan . That is the shock: swans are black as well .
    \item The same reasoning runs through everyday life . You meet a few people from a minority community for the first time, observe how they behave, pray and eat, and then make a generalised statement about all of them . Social media is nothing but a reflection of our inductive reasoning .
\end{itemize}
\end{ExampleBox}

\begin{CriticalPoint}
\begin{itemize}
    \item \textbf{Dictated problem:} The biggest problem with inductive reasoning is \textbf{hasty generalisation} .
    \item This was articulated by \textbf{David Hume}---a scholar of the Scottish Enlightenment, who talked about scepticism, doubt and empiricism .
\end{itemize}
\end{CriticalPoint}

\paragraph{The dictated characteristics of inductive reasoning}
\begin{itemize}
    \item This method of reasoning involves moving away from the \textbf{specific towards the general}, and finally a theory is developed .
    \item Inductive reasoning involves: $\text{Data collection} \longrightarrow \text{Classification} \longrightarrow \text{Theorisation}$ .
    \item Which means there is no method involved to test inductive reasoning . Generalisation is developed without testing, purely on the basis of observation .
    \item Before saying ``all swans are white'', the statement was never tested . Data was simply collected---four white swans one day, ten the next, seventeen the day after---and then generalised .
    \item If experiments are done at all, they are done to prove what you already think is right . Having assumed that all swans are white, you collect data only to prove yourself right and avoid anything like a black swan, because it invalidates your theory .
\end{itemize}

\begin{ExampleBox}
\textbf{The prelims example:} In the examination you see a question whose answer you do not know and you mark option C . You come home and start collecting data to prove that C is right . If someone says B is correct, you produce multiple sources to show him that he is wrong and you are right . Is this scientific? No . You are collecting data to prove yourself right while avoiding anything that could prove you wrong .
\end{ExampleBox}

\begin{CriticalPoint}
That is why inductive reasoning, which was initially very popular, has today been discarded: people tend to collect data to prove their theory right and avoid the other picture, so you do not get the absolute or objective reality .
\end{CriticalPoint}

\begin{QuickRevision}
\begin{itemize}
    \item Induction moves from the specific to the general; theory is built at the end .
    \item Sequence: data collection $\rightarrow$ classification $\rightarrow$ theorisation; no testing involved .
    \item Core defect: hasty generalisation---David Hume (Scottish Enlightenment; scepticism, doubt, empiricism) .
    \item Black swan, minority stereotyping, option-C data collection---all illustrations of the same error .
\end{itemize}
\end{QuickRevision}

\subsection{Deductive (Hypothetico-Deductive) Reasoning}
The first example of the class---the car, with its prediction, experimentation and testing---was deductive reasoning all along .
\begin{itemize}
    \item Write the heading as deductive or hypothetico-deductive; both are one and the same . In some classical books the hypothetico-deductive method is described as being coined by \textbf{Karl Popper}, and other books may use other names; for our purposes the two terms are treated as identical to avoid confusion .
    \item \textbf{Dictated:} This method of reasoning, unlike the inductive method, is based on testing . The claim could be right or could be wrong, but it has to be tested .
    \item Having seen ten or twenty swans does not entitle us to be certain that all swans are white . We have to test---to go to multiple countries---instead of limiting ourselves to Europe and declaring from there that all swans are white .
\end{itemize}

\begin{CriticalPoint}
\begin{itemize}
    \item This is exactly the kind of generalisation Europeans made about non-European people while sitting in Europe .
    \item \textbf{Armchair philosophers:} Those who sit in their chair and theorise about the world without having gone anywhere outside their own location . This is inductive reasoning, nothing else .
\end{itemize}
\end{CriticalPoint}

\begin{QuickRevision}
\begin{itemize}
    \item Deductive = hypothetico-deductive (in some books associated with Karl Popper) .
    \item Unlike induction, it is based on testing a hypothesis you already hold .
    \item Armchair philosophy = generalising about the world without leaving your location = induction .
\end{itemize}
\end{QuickRevision}

\subsection{Verifiability, Falsifiability and the Provisional Nature of Science}

\begin{KeyConcept}
\textbf{Dictated conclusion:} \textit{In conclusion, the inductive method is based on the verifiability principle and the deductive (hypothetico-deductive) method is based on the falsifiability principle.} 
\begin{itemize}
    \item \textbf{Verifiability} is only a more flowery word for the same swan story: 10 white swans, then 20, then 30, then 100; the generalisation ``all swans are white'' follows; and thereafter, wherever I go, I try only to prove this statement right .
    \item Having concluded, I keep verifying and am not ready to refute; I avoid any evidence that contradicts my theory . No scientist likes to be contradicted---every scientist would want a theory that explains everything around them .
    \item However, this should be avoided, and that is why people now follow deductive reasoning more than inductive reasoning . Even in the natural sciences we rely more on the test .
\end{itemize}
\end{KeyConcept}

\begin{ExampleBox}
\textbf{Einstein's theory of relativity:} Its essence, as stated in class: \textit{nothing in the world travels faster than light} . Is the statement testable? Yes . If you find something travelling faster than light, the statement is rejected---but that too shows it was tested . If you find nothing faster, the statement remains valid, again because it was tested . That is why the theory of relativity is scientific---because it is testable .
\end{ExampleBox}

\begin{CriticalPoint}
\begin{itemize}
    \item Something being scientific does not mean it is absolute truth . It only implies that it is testable .
    \item Therefore \textbf{scientific knowledge is always provisional---temporary} .
\end{itemize}
\end{CriticalPoint}

\begin{DataFact}
\begin{itemize}
    \item In 2011, scientists from Italy reported sub-atomic particles travelling faster than light (quantum physics) . They could not prove it .
    \item The failure to prove is a separate matter; what the episode shows is that the statement is testable .
\end{itemize}
\end{DataFact}

\paragraph{The dictated chain}
\begin{itemize}
    \item A statement that is scientific only implies that it is testable .
    \item If something is testable, it means it can be falsified as well .
    \item And if something can be falsified, that means that knowledge is provisional---the so-called scientific knowledge is provisional, temporary .
\end{itemize}

\paragraph{Newton, Einstein and falsification}
\begin{itemize}
    \item Is Newton's law of gravitation testable? Of course---and therefore scientific .
    \item The problem is that Einstein believed Newtonian physics cannot explain everything in nature . Newton's theory could not explain why the orbit around Mercury is formed as it is; there are many things it cannot explain .
    \item At that point of time Newtonian physics was believed to be absolutely scientific and was the dominant paradigm in the sciences, able to explain everything . Newton was virtually a god figure: if Newton had said something, it was absolutely right---until Einstein came in and rejected many of his theories .
    \item Why could Einstein reject them? Because those theories were testable---that is, they could be falsified .
\end{itemize}

\begin{ExampleBox}
\textbf{Space-time sheet:} To Newton, gravitation is a force . Einstein said clearly that gravity is not a force . Take a sheet of paper as space and drop a ball on it: the sheet folds into that shape . To Newton this illustrates gravity; to Einstein it is space compression---the warping of space---and no such force as gravity exists . Einstein's own theory of relativity is likewise regarded as falsifiable---it can be falsified if we find something travelling faster than light . People have found candidate phenomena; they have not been able to prove them .
\end{ExampleBox}

\begin{DataFact}
\begin{itemize}
    \item There was a time when we had knowledge of nine planets; that was not absolute . Today we study eight, and the ninth was found not to be a planet at all .
    \item For any body to be called a planet, certain criteria must be met; if they are not met, it is not a planet . Science too grows and matures over time .
\end{itemize}
\end{DataFact}

\begin{KeyConcept}
\textbf{Prediction, observation and experimentation} are the basic postulates---the basic features---of scientific method . A method is called scientific only if through it you can test something: that is, falsify, predict, or experiment in order to test . Whether or not you wear a white suit, and whether you are in a car or in a lab, does not matter .
\end{KeyConcept}

\begin{QuickRevision}
\begin{itemize}
    \item Induction $\rightarrow$ verifiability; deduction $\rightarrow$ falsifiability .
    \item Scientific $\neq$ true . Scientific = testable, hence falsifiable, hence provisional .
    \item Relativity is scientific because testable; Newton could be rejected because his theory was falsifiable .
    \item Nine planets $\rightarrow$ eight planets: science matures; knowledge is temporary .
\end{itemize}
\end{QuickRevision}

\subsection{Doubts Discussed in Session 1}

\paragraph{1. Can you explain social psychology with reference to Unit 1?}
\begin{itemize}
    \item Social psychology is more about the study of group behaviour .
    \item Conventionally, sociology is about structures---family, religion---and their impact on individuals . Psychology is about the individual only: why you are a criminal, what is wrong with you; not religion, economy or globalisation .
    \item Social psychology asks why a certain group behaves as it does, or how an individual is influenced by his peers and classmates .
\end{itemize}

\begin{ExampleBox}
On a messaging application there are multiple sociology channels and groups where people discuss answer writing and share answers for review . You tag someone whose identity you do not really know and ask him to review your answer, because he has presented himself as someone reliable . That group then influences your behaviour: you see people discussing economy or inequality or some other book, and you also want that book, those notes, or that person's help . The ``group'' need not be an online group---your family and your classroom influence your behaviour in the same way . Social psychology therefore falls directly in between the larger structural study of sociology and the micro study of psychology .
\end{ExampleBox}

\paragraph{2. Is the theory of the presence of life on other planets scientific?}
\begin{itemize}
    \item First, a related question: India is habitable; Pakistan is habitable; China is habitable; Europe is habitable---and from these four cases I conclude that the Earth is habitable . What kind of reasoning is this? It is inductive, because I move from the specific to the general .
    \item The generalisation is unsafe: Antarctica is not habitable; deserts and other natural landscapes do not offer conditions conducive to survival (the Sentinelese Island was mentioned and set aside as a controversial case) . So the Earth in itself is not habitable, though for the sake of argument we accept that it is .
    \item Next I begin looking for Earth-like planets in the outer world . A newspaper reports that a planet very similar to Earth, many light years away, appears to have Earth-like conditions . Is that report scientific? The question is simply: can it be tested? It can---one could go there and see whether one can live there . There must, of course, be criteria to test something by .
\end{itemize}

\begin{ExampleBox}
\textbf{Life on Mars:} If it is claimed that Mars is suitable for us, can it be tested? Of course . You send a rover; experiments are done with animals sent up, to see how they react---whether they lack oxygen, whether water can be found . A machine on the rover collects samples of Martian soil; that data is processed in laboratories on Earth to see whether elements such as nitrogen---the basic element required for human existence, the building block of protein---are present; oxygen and hydrogen come later . If some elements are missing, further rover missions follow, until all the elements required for human life are verified or ruled out . Everything depends on rigorous experimentation . The answer to the doubt: yes---the only measure of whether something is scientific is whether it can be falsified, that is, tested .
\end{ExampleBox}

\paragraph{3. Why should we define social action in an answer on sociology and common sense?}
Because when you study social action you are studying the common sense of the individual . Every individual attaches meaning to his or her action, and even to institutions .

\begin{ExampleBox}
Family is undoubtedly socialising you---but the individual also uses the family to gain benefits in the outer circle . If you belong to a family of administrators or bureaucrats, you are regarded as superior among your friends . How we use family is therefore equally important . Just as it is important for a structural functionalist to analyse how family contributes to the system, it is important for the sociologist of social action to understand how the individual uses family . \textit{``If you are the son of a king, an emperor, a bureaucrat or a prime minister, you will act like this''}---that is common sense . But we are studying this common sense too . That is why social action brings sociology closer to common sense .
\end{ExampleBox}

\paragraph{4. Are hypothesis and reasoning the same or different?}
Hypothesis is a part of reasoning . You reasoned out multiple causes behind the failure of your car; all those reasons or causes were hypotheses only, none of them yet tested or proven right . You try to prove one of them right; if it is not, you come up with another---this may be the cause---and so it goes on .

\paragraph{5. What is romantic ethics?}
Some words in the handouts are not yet covered; that is natural .

\begin{ExampleBox}
You have just bought a lipstick . The next day you look at an online cosmetics store and come across shades better than the one you purchased, and you now want those too . It is not the difference that makes you buy---it is the novelty . The same with gaming consoles: if I have one generation and the next is released, I want it; and then the one after that . If you have one smartphone model you want the next .
\end{ExampleBox}

\begin{KeyConcept}
\begin{itemize}
    \item \textbf{Romanticism} = falling in love with something new; it is human nature to want the latest gadgets and products .
    \item \textbf{Romantic ethics} creates the spirit of consumerism: we seek novel experiences and products .
    \item Economists study only that GDP has increased; what went into that increase---our fetish for something new---cannot be mathematised, graphed or made statistical . You cannot compute desire; there is no percentage of desire .
\end{itemize}
\end{KeyConcept}

\paragraph{6. Can we say the inductive method creates new theories and the deductive method reforms old theories?}
Yes---there is no doubt about it . The ability to test theories is purely deductive; the ability to develop new theories is purely inductive---you observe a few things and generalise, a theory built purely on observation .

\begin{ExampleBox}
Why did you think in the first place that the battery could be the reason? Because since childhood you had seen two or three cars that would not start because of the battery . From those specific instances you developed a general expectation---that is inductive reasoning . The first hypothesis was therefore itself produced by induction, and was then tested through the deductive process .
\end{ExampleBox}

\paragraph{7. Can we say inductive knowledge is closer to common sense?}
Let us not say anything like that for the time being . You will automatically get the answer by the end of this unit .

\paragraph{8. Does that not mean that the values and morals of a person will affect theory and study?}
Yes---values and morals affect our theories, especially in the social world . In the natural world theories are believed highly objective and unaffected by scientists' values---but they are affected too: experiences shape results, and worldview and social location influence the outcome . You will realise this fully in the critique of scientific method: those who employ scientific method and call themselves scientists are not always objective .

\begin{ExampleBox}
\textbf{Eugenics:} Described in class as a science which deals with the study of race . It was used to justify the superiority of the white race: the subject appears scientific, but science was being used by the researcher to justify superiority over so-called uncivilised people . Having declared them uncivilised, a discipline was invented to justify the declaration . Science cannot always be objective .
\end{ExampleBox}

\begin{QuickRevision}
\begin{itemize}
    \item Social psychology = group behaviour, between sociology and psychology .
    \item Life on other planets is scientific because testable (rover, soil samples, nitrogen) .
    \item Social action = the individual's own common sense $\rightarrow$ sociology closer to common sense .
    \item Hypothesis is part of reasoning; induction builds theory, deduction tests it .
    \item Romantic ethics $\rightarrow$ novelty $\rightarrow$ consumerism; desire cannot be computed .
    \item Scientists' values and social location shape results---eugenics .
\end{itemize}
\end{QuickRevision}

\sectiondivider

\section{Features of Science, Norms of Science and the Scientific Method}

\subsection{Opening Question: Is Mathematics a Philosophy?}
The session opened with a deliberately simple question: is mathematics a philosophy?  To those who answered no, a basic counter-question was put---if mathematics is not philosophy, then why do we write PhD in mathematics? What is the full form of PhD? Doctor of Philosophy . On that reasoning, mathematics should be philosophy .
\begin{itemize}
    \item Philosophy is the mother of all sciences . Sociology therefore cannot pose any challenge to philosophy; sociology is too young to resist or protest against it . Philosophy has existed since human life came into existence .
    \item Philosophy is not a subject---not a three-year BA programme . It is a set of knowledge: a systematic set of knowledge, not common sense, having its own unique method, its own history of origin in the form of a discipline, and its own contributions .
    \item Those contributions have enhanced and enriched the way we look at the world around us---all of them made by philosophers .
\end{itemize}

\begin{KeyConcept}
\begin{itemize}
    \item Mathematics is a philosophy in this sense: a systematic set of knowledge with its own method and history .
    \item Philosophy = mother of all sciences; sociology is too young a discipline to be set against it .
\end{itemize}
\end{KeyConcept}

\begin{QuickRevision}
\begin{itemize}
    \item PhD = Doctor of Philosophy---the argument used to show mathematics is philosophy .
    \item Philosophy is a systematic set of knowledge, not a course; it is the mother of all sciences .
\end{itemize}
\end{QuickRevision}

\subsection{How the Question Is Framed: The Problem With Survey Data}
The question was not asked neutrally . It was asked with a particular expression and emphasis---is mathematics a philosophy?---and that manner of asking may itself have pushed some students towards answering no .

\begin{CriticalPoint}
\begin{itemize}
    \item This is not a problem with one instance . It is the problem with surveys: the way questions are framed, and the way they are posed by the person asking, also influences the data collected .
    \item Had the data been gathered on ``is mathematics philosophy?'' while asking it in that dismissive tone, the maximum responses would have been ``no''---and that would not mean mathematics is not philosophy .
    \item So merely collecting data does not mean the theory is valid; merely observing and collecting data does not mean you will reach the true conclusion .
\end{itemize}
\end{CriticalPoint}

\begin{QuickRevision}
\begin{itemize}
    \item Question framing and the questioner's manner influence survey data .
    \item Data collection by itself guarantees neither validity nor truth .
\end{itemize}
\end{QuickRevision}

\subsection{What Is Science? The Five Features}
Science, as said in the previous session, is a way of life . But for the answer sheet the definition to be written is: \textbf{science is a body of knowledge characterised by certain features} . Five features were dictated:
\begin{enumerate}
    \item \textbf{Testability:} Anything dealt with in science should be testable . If it is not testable, it is not scientific . When you test something, you are actually following scientific method .
    \item \textbf{Empirical:} Empiricism means observation with your senses . Only if a thing comes to your observation through your senses can you call it empirical .
    \begin{itemize}
        \item There are many studies where direct observation is impossible---you cannot observe a chemical reaction at the molecular level with the naked eye . But if you observe it under a microscope, that is still empirical observation, because you are observing through your own eyes with the help of a scientific instrument that lets you observe more closely .
    \end{itemize}
    \item \textbf{Objectivity:} The values, biases and opinions of the scientist conducting the research should be kept aside . The research should be completely value neutral or value free---the two terms need not be worried about at this stage .
    \begin{itemize}
        \item Objectivity = the act of observing without biases, opinions or values . Empiricism means observation; that observation must be objective---not marred or affected by the researcher's values, biases or opinions . Only then can we reach the objective truth or objective finding .
    \end{itemize}
    \item \textbf{Cumulative:} Cumulative means the knowledge has added to something already known . Just as we speak of accumulation---gathering, building up---cumulative knowledge builds on already existing knowledge . You would not have arrived at this knowledge had there not been a knowledge prior to it .
    \item \textbf{Theoretical:} Whatever data you have collected through observation or empiricism, and whatever you have analysed from cumulative knowledge, is finally made into a theory that can explain the world or the laws which govern the natural world .
    \begin{itemize}
        \item That is the role of theory: we take the help of theories to explain why something happens in nature . Theory is therefore obviously a feature of science .
    \end{itemize}
\end{enumerate}

\begin{ExampleBox}
It is easy for us today to talk about the structure of the universe---whether the sun or the earth is at the centre . But Copernicus could not have developed the heliocentric model if there had been no prior model placing something---the earth---at the centre . Only by rejecting it could he advance his own knowledge . That is why a scientific theory always comes as the result of the refutation of an already existing theory . Science builds up .
\end{ExampleBox}

\begin{KeyConcept}
\begin{itemize}
    \item \textbf{Five features of science to be remembered:} Testability, empiricism, objectivity, cumulative, theoretical .
    \item \textbf{Definition to write:} Science is a body of knowledge characterised by testability, empiricism, objectivity, cumulativeness and theory .
\end{itemize}
\end{KeyConcept}

\begin{QuickRevision}
\begin{itemize}
    \item Testability---untestable means unscientific .
    \item Empirical---observation through the senses; instrument-aided observation still counts .
    \item Objectivity---observing without biases, opinions or values (value neutral / value free) .
    \item Cumulative---builds on and refutes prior knowledge (Copernicus needed geocentrism to reject) .
    \item Theoretical---data finally becomes theory explaining the world .
\end{itemize}
\end{QuickRevision}

\subsection{Testability Examined: Hypnosis and Aliens}

\begin{ExampleBox}
\textbf{The hypnosis episode:} A friend tried hypnosis on me . Knowing in advance what he intended, I acted as if I had been hypnotised---so he could not really hypnotise me . After an hour-long session he declared that hypnosis works, and told a third person what was supposedly going on in my mind, what kind of man I am, what my ambitions are---everything ``hidden in the unconscious layer'' and revealed through hypnosis . Can that ever be tested? There is no test available to check whether what he observed through hypnosis was actually going on in my unconscious mind . It cannot be analysed in a graph or given any visual appearance that would prove it . If I cannot test it, I can never falsify it; if I cannot falsify it, I cannot call hypnosis scientific---even though it is defended by pointing to Sigmund Freud and psychoanalysis, and to psychology as a scientific discipline .
\end{ExampleBox}

\begin{ExampleBox}
\textbf{``Aliens exist in the outer world'':} Such news is sometimes carried even in scientific journals . How would you test it? Just as a mathematician begins with an assumption---let us assume $x = 1$---the assumption must then be testable; in mathematics there are formulas to test it . Here there are none . Suppose we say: let us go and explore that place . Till what time will you explore it? From a very high altitude you photograph an object and say it looks like a human-like figure---but first descend to that place . And who told you that aliens look like human beings? That is one more assumption . Because something looks human, we call it an alien of the outer world . Our observation is being influenced by assumptions, biases and opinions . If it cannot be tested, it is not scientific . As simple as that .
\end{ExampleBox}

\begin{CriticalPoint}
A question was left deliberately open at this stage: is yoga a science? Students were told their answer will change once the features of science are properly applied, and that the body of knowledge must be completely objective---not called scientific out of emotion merely because yoga is a gift of India to the world .
\end{CriticalPoint}

\begin{QuickRevision}
\begin{itemize}
    \item Hypnosis: claims about unconscious content cannot be tested $\rightarrow$ cannot be falsified $\rightarrow$ not scientific .
    \item ``Aliens exist'' is untestable and rests on the further assumption that aliens look human .
    \item Objectivity forbids calling something scientific out of national or emotional attachment .
\end{itemize}
\end{QuickRevision}

\subsection{Norms Governing Science --- Merton's Kudos Principle}
After the features of science come the norms governing science . Four norms were dictated:
\begin{enumerate}
    \item \textbf{Communalism:} This word does not carry its everyday Indian meaning here . Communalism means that whatever scientific contributions or discoveries scientists come up with belong to the world; nobody should restrict anyone's access to a contribution that advances humanity .
    \begin{itemize}
        \item \textit{Dictated definition:} Communalism is a norm in sciences according to which any scientific invention or discovery should be shared with the larger community, that is, the world .
    \end{itemize}
    \item \textbf{Universalism:} Whether you are a black scientist, a female scientist, a Dalit scientist, an American or African scientist, a disabled scientist, a lesbian or gay scientist---at the end of the day you are a scientist only . Science does not discriminate; if you have discovered something you will be celebrated, irrespective of your social location .
    \begin{itemize}
        \item \textit{Dictated definition:} Universalism is defined as acknowledgement of the scientific contribution irrespective of the social location of the scientists . It simply means that the age, colour, caste, creed, sex, sexuality, religion or nationality of the scientist do not matter .
    \end{itemize}
    \item \textbf{Disinterestedness:} If I am a scientist and I have invented something, it should not be used for my personal or private gain .
    \begin{itemize}
        \item Research is possible because of funds given by the state---for instance to a scientist in DRDO or CSIR . If you come up with an invention or discovery, you should be thankful to the state, and whatever you have developed cannot be used for personal gain; it has to be shared with everyone .
        \item Disinterested therefore means you are not interested in using the thing to pursue your private interest---you are altruist, not interested in your own self but in the welfare of the masses . It is very similar to communalism .
    \end{itemize}
    \item \textbf{Organised Scepticism:} Scientists as a community should be ready for criticism---for someone to criticise their finding or find loopholes in it---without feeling angry . They maintain restraint and say: thank you, I will work on that .
    \begin{itemize}
        \item The reasoning behind the norm: we are not working for ourselves, we are working for the world; so criticism is largely for the welfare of communities and societies and should not be taken to heart .
        \item When Einstein criticised Newton, Newton did not feel angry . When scientists criticise each other they do not take it to heart---whereas when social scientists criticise each other it is another game altogether .
    \end{itemize}
\end{enumerate}

\begin{ExampleBox}
\begin{itemize}
    \item \textbf{Communalism example:} If I am Thomas Alva Edison and I have invented the electric bulb, that bulb is everyone's property, because it contributes to the advancement of humanity . If I am Alexander Graham Bell and I invented the telephone, my invention must be shared with the world community .
    \item \textbf{Universalism example (DNA structure):} Who discovered the structure of DNA? The immediate answer given by everyone is Watson and Crick, because they received the Nobel Prize . It was in fact Rosalind Franklin, a female scientist . There was a dispute going on between the three; scientists too are human beings; and at the end of the day the award went to the other two . The very speed of the classroom answer shows that discrimination exists in the scientific community .
    \item \textbf{Universalism example (Einstein):} Albert Einstein---many of his claims were initially refused or rejected by British and French scientists . Einstein was a German, and a war was going on between England and Germany; the French and the English had decided to boycott Germans, so even a great innovation would be rejected purely because it came from Germany . And if science does not discriminate, how do you explain the overpopulation of men and underpopulation of women in STEM fields---science, technology, engineering, mathematics? Science often does discriminate .
\end{itemize}
\end{ExampleBox}

\begin{CriticalPoint}
Is it a law or only a norm? It is a norm . Norms can be broken---and so can laws---but breaking a norm does not send you to jail, while breaking a law does . Example given: Bill Gates decided not to share the vaccine formula with the underdeveloped world, the third-world nations including India . The norm of communalism was thereby broken---and that is precisely why it became news .
\end{CriticalPoint}

\begin{KeyConcept}
\begin{itemize}
    \item These four norms were formulated by \textbf{Robert K. Merton}, also called a sociologist of science, who is in the syllabus .
    \item They are collectively known as the \textbf{CUDOS / KUDOS principle}---the acronym formed from the initials of the four norms (Communalism, Universalism, Disinterestedness, Organised Scepticism) .
    \item The norms are valid for both natural science and social science; both are science only .
\end{itemize}
\end{KeyConcept}

\vspace{6pt}
\noindent
\begin{tabularx}{\textwidth}{p{3.2cm} Y Y}
\toprule
\rowcolor{AccentDark}
\thead{Norm} & \thead{What it requires} & \thead{How it is broken} \\
\midrule
\textbf{Communalism} & Discoveries belong to the world community and must be shared & Withholding a vaccine formula from poorer nations  \\
\textbf{Universalism} & Contribution acknowledged irrespective of social location---age, colour, caste, creed, sex, sexuality, religion, nationality & Credit for the DNA structure denied to Franklin; rejection of Einstein for being German; gender gap in STEM  \\
\textbf{Disinterestedness} & No private or personal gain from state-funded research; altruism & Using an invention to pursue private interest  \\
\textbf{Organised Scepticism} & Readiness to be criticised, with restraint & Taking criticism to heart, as often happens among social scientists  \\
\bottomrule
\end{tabularx}
\vspace{6pt}

\begin{QuickRevision}
\begin{itemize}
    \item Four norms---Communalism, Universalism, Disinterestedness, Organised Scepticism = KUDOS principle, given by Robert Merton .
    \item They are norms, not laws: breakable without legal penalty .
    \item Each was illustrated by a case of its violation .
\end{itemize}
\end{QuickRevision}

\subsection{The Scientific Method: From Observation to Publication}
The car example was revisited to formalise the steps: I went out of my house intending a casual round; I saw that my car was not working . That was observation .
\begin{enumerate}
    \item \textbf{Observation:} Car is not working .
    \item \textbf{Hypothesis:} Many possible reasons came to mind; I believed the battery was the issue .
    \item \textbf{Prediction:} If I replace the battery, everything will be fine .
    \item \textbf{Experimentation:} I obtained a battery from my brother and replaced the old one .
    \item \textbf{Confirmation:} The car now works; my hypothesis was right . Had it not worked, the hypothesis would have been nullified (null hypothesis) and I would have moved to the alternative hypothesis---perhaps the engine---and tested it by the same steps .
\end{enumerate}

\begin{CriticalPoint}
\begin{itemize}
    \item Your finding can be confirmed or disconfirmed . Would you cry if your hypothesis were disconfirmed? No---you should be glad that you now know this is not the cause .
    \item Every time a hypothesis is discarded you are moving closer to the truth . Celebrate validation, but do not grieve over rejection; have patience instead .
\end{itemize}
\end{CriticalPoint}

\begin{KeyConcept}
\textbf{Dictated definition:} \textit{Scientific method consists of systematic steps to collect and analyse data in an objective and theoretical fashion, to conclude the finding or to reach the truth.} 
These steps are followed not only by natural scientists but by everyone who claims a scientific approach and belongs to the scientific community---sociologist, psychologist, economist, geographer, mathematician, chemist or physicist . The steps remain similar for all of them .
\end{KeyConcept}

\paragraph{Step 1 --- Observe and define the problem}
First you observe the problem, which automatically implies defining it . Otherwise there would be no reason to follow the method at all .

\begin{ExampleBox}
In the time of Copernicus the central problem was: what exists at the centre of the universe---the earth or the sun?  That was the problem he identified and defined . Put formally: Copernicus identified the problem with the geocentric theory .
\end{ExampleBox}

\paragraph{Step 2 --- Review the literature}
\begin{itemize}
    \item \textit{Dictated:} This step involves investigation of the already existing literature in order to find the mistakes or the research gap .
    \item Copernicus, before developing his own thesis, went through the work of Ptolemy, a great astronomer, and of Aristotle---reading everyone else to see what they had to say on the model and why .
    \item When your car was not working you also went through manuals and handouts, or looked up a video . Literature does not mean text only: you are reviewing the work already done by others on the same issue, and how they tackled it .
\end{itemize}

\paragraph{Framing the research question precisely}
\begin{itemize}
    \item For higher studies you will be asked what topic you will work on, and the topic must be defined precisely and specifically---not vague or open-ended .
    \item You cannot say ``I want to study why India has poverty''---an individual cannot study the whole of India . Better: why poverty exists in a particular district, or a particular village in a state---that is feasible .
    \item More precise still: \textit{``Why does poverty exist among the Dalits of village A?''}  Open-ended questions are not encouraged; a supervisor will accept you only when the research question is narrow and clear .
\end{itemize}

\begin{ExampleBox}
You go to the university library, and search on Google Scholar, which tells you what has already been researched and what scope is left . Thousands of pages appear . Going through them, you find that studies exist on poverty among women, or among children, or among a particular caste group---but nobody has yet studied why poverty exists among the Dalits of village A . That is your \textbf{research gap}, and this is what you can now research on .
\end{ExampleBox}

\begin{CriticalPoint}
An aside on the research culture: people often do research in order to avail their fellowship . Scientific spirit is not yet there . A PhD that can be completed in three years abroad takes a minimum of five years in India, because the Junior Research Fellowship (JRF) is valid for five years, along with accommodation and other resources .
\end{CriticalPoint}

\paragraph{Step 3 --- Formulation of hypothesis}
\begin{itemize}
    \item Hypothesis, as already known, is any assumption that is yet to be tested . \textit{Dictated:} Having gone through tons of literature, the researcher develops one hypothesis .
    \item After reading articles and journals by different scholars---some published in \textit{Economic and Political Weekly} (EPW), some with a ministry of labour---you will have developed certain assumptions . Every mathematician or scientist starts with an assumption; so do you .
    \item The hypothesis framed here: \textit{``Poverty exists among the Dalits of village A because of lack of education.''}  Is it right or wrong? It is yet to be tested---exactly as with the battery .
\end{itemize}

\begin{CriticalPoint}
In the social sciences, and especially in sociology, we do not solve the problem---we study the causes of the problem . Initial sociologists had the ambition to solve problems . Today we do not .
\end{CriticalPoint}

\paragraph{Step 4 --- Data collection and analysis}
Data is collected through two routes :
\begin{itemize}
    \item \textbf{Primary source:} I myself go to the village and observe empirically, through my own senses . This is field study .
    \item \textbf{Secondary source:} Someone else's work: what was already done at the review-of-literature stage . Newspaper reports, journals, science journals, EPW . This is book study or armchair study; the two terms are the same . Note that those authors may themselves have relied on secondary literature .
\end{itemize}

\begin{ExampleBox}
\textbf{Going to the field:}
\begin{itemize}
    \item You take a flight to the state capital, then a train or bus, and reach the village . You stay at least a year, rent a house, and maintain cordial relations without letting people know you are an outsider from the city---because a close-knit village would not share personal or intimate information with an outsider .
    \item Asking directions: you cannot ask where the ``depressed'' or ``oppressed classes'' live; you have to speak as people in the field speak, and ask where the Dalit residential colony is . You are warned not to go there, that you will be ``polluted'' .
    \item You knock at a Dalit household and ask for a glass of water . The question itself carries a normative order you may not be aware of; the door is shut . Days pass . You begin by befriending the children, giving them small change for sweets, without even knowing their language .
    \item The problem is multi-fold: you do not know their language, their norms, or anything about them---you only know that you must finish the research and submit it .
    \item In the field, data is collected through \textbf{interviews, surveys, a questionnaire, or participant observation} . (These recur in Unit 3 as well.) 
    \item To answer why Dalits are poor in this village, you first have to know who the Dalits are, what their education is, what the composition of the Dalit population is, and how many there are---and Dalits are not a homogeneous category; Dalit women and their issues are also there .
\end{itemize}
\end{ExampleBox}

\paragraph{Classification and pattern}
\begin{itemize}
    \item After collection comes classification, to make the data simpler---exactly as you make a flow chart or a table when you understand a topic .
    \item Classification is how we make sense: unless you classify binary opposites such as good and evil, you cannot differentiate between them; each exists because the other does .
    \item Having interviewed, surveyed and questioned, say, 1000 Dalit people, you classify by gender---for example 700 male and 300 female---and then further, by education: of 700, 600 up to class five and 100 with no education . Classification is already the beginning of analysis .
    \item Sociologists always look out for a pattern in the classified data . The simplified data used in class:
\end{itemize}

\vspace{6pt}
\noindent
\begin{tabularx}{\textwidth}{p{2.5cm} p{4cm} p{2.2cm} Y}
\toprule
\rowcolor{AccentDark}
\thead{Category} & \thead{Education (Max.)} & \thead{Sexuality} & \thead{Class / Economic Status} \\
\midrule
\textbf{Dalit male} & Under class five: 500 males; 500 who studied till class twelve & All recorded as straight & Largely middle class---sanitation and similar occupations; better placed than those with no schooling  \\
\textbf{Dalit female} & Under class five: 100 females; largely no education & All recorded as straight & Lower class; mainly working in other households to earn their bread---domestic help is largely drawn from the Dalit caste  \\
\bottomrule
\end{tabularx}
\vspace{6pt}

\begin{KeyConcept}
Going back to the research question after analysis, the answer is not only lack of education . The answer is also \textbf{patriarchy} .
\begin{itemize}
    \item Dalit males did receive at least some education; Dalit females did not . So behind the poverty of Dalit women lies lack of education---but behind that lack of education lies Dalit patriarchy .
    \item In the pattern that emerges, the Dalit woman is \textbf{doubly oppressed}---she is Dalit and a woman . Dalit is generally considered the depressed class by Ambedkar; the Dalit woman is further depressed than that depressed class . This may explain why women are placed at the worst end of the spectrum among Dalits .
\end{itemize}
\end{KeyConcept}

\paragraph{Step 5 --- Publication of the findings}
\begin{itemize}
    \item The findings are finally submitted and published: either publishers approach you, or you submit your research to a publication house .
    \item If your research is relevant to their journal it will be published; if it is not scientific enough---not based on any data, only common sense---it will never be accepted . You must show that it is rigorous, extensive research with data to back the theory and the findings .
\end{itemize}

\begin{ExampleBox}
An article on the platform economy---a study of app-based drivers, conducted with a colleague from a school of economics---was published by a national labour institute under the Ministry of Labour, because it was based on data and scientific research and fitted their scientific paradigm .
\end{ExampleBox}

\begin{QuickRevision}
\begin{itemize}
    \item Method as practised: observation $\rightarrow$ hypothesis $\rightarrow$ prediction $\rightarrow$ experimentation $\rightarrow$ confirmation/disconfirmation .
    \item Method as formalised: define the problem $\rightarrow$ review the literature (find the research gap) $\rightarrow$ formulate hypothesis $\rightarrow$ collect and analyse data $\rightarrow$ publish findings .
    \item Definition: systematic steps to collect and analyse data in an objective and theoretical fashion to reach the truth .
    \item Primary source = field study; secondary source = book/armchair study .
    \item Worked case: poverty among Dalits of village A $\rightarrow$ hypothesis of lack of education $\rightarrow$ finding of patriarchy and the double oppression of the Dalit woman .
\end{itemize}
\end{QuickRevision}

\subsection{Critique of Scientific Method I --- Karl Popper and Demarcation Principle}
The last component of sub-unit 2A is the critique of scientific method . We have studied science and the scientific method, which natural and social scientists both practise . The question now is: is scientific method scientific? This is the biggest problem, especially in the social sciences .
\begin{itemize}
    \item \textbf{Karl Popper}, one of the greatest philosophers of science, gave the \textbf{demarcation principle}---the principle which demarcates science from pseudoscience .
\end{itemize}

\begin{CriticalPoint}
\begin{itemize}
    \item What is pseudoscience? The term appears in the newspapers---for instance the Indian Medical Association calling yoga a pseudoscience, that is, not science .
    \item \textbf{An answer-writing caution given here:} Avoid political personalities in your answers . Write \textit{``the Indian Medical Association called yoga a pseudoscience''}, not a sentence naming political leaders or their ventures . Write in a scientific or administrative fashion .
\end{itemize}
\end{CriticalPoint}

\paragraph{Popper's criteria}
One simple criterion follows from the five features already learnt: testability, empiricism, objectivity, theoretical, cumulative---if something does not have these features, it is pseudoscience . Popper made it simpler still . Science, according to Popper, has two important elements :
\begin{enumerate}
    \item We follow the \textbf{falsifiability approach} .
    \item We follow \textbf{deductive reasoning} .
\end{enumerate}
True science is the one which follows falsifiability and deductive reasoning, because both imply testability . If you can falsify something, you are testing it; and science involves rigorous testing .

\begin{ExampleBox}
\textbf{Kekulé and the benzene ring:}
\begin{itemize}
    \item Kekulé discovered the structure of benzene in a dream---he saw snakes running around a campfire, woke up, and drew exactly what he had seen; the cyclical ring, with the snakes standing for the double bonds .
    \item But the dream by itself did not settle the matter: it had to be tested rigorously in the laboratory . He conducted experiments and concluded that without the double bonds the reaction would not have proceeded---hence the double bonds exist, and how many .
    \item Would Popper call it scientific? Definitely .
    \item And is it inductive or deductive? Deductive---because you already have a hypothesis and then test whether it is right or wrong . In induction you develop the hypothesis after observations (10, 20, 30 white swans and then the generalisation); there was no test there . In deduction you always test something you saw---in a dream, or heard from someone .
\end{itemize}
\end{ExampleBox}

\begin{ExampleBox}
\textbf{The dream of a deity:}
\begin{itemize}
    \item Someone says a deity appeared in his dream and instructed him to build a temple at a particular place, and calls his friends to begin . He cannot prove that the vision occurred . If you cannot prove something, how can it be scientific? 
    \item And this is exactly how a conflict between science and religion arises . Science runs on empiricism, and doubt is the central feature of science . When you doubt religious people they start hating you; when you doubt in science, we celebrate---because doubt is the step one must take before advancing a new theory .
    \item Pseudoscience, by contrast, does not falsify anything---it only verifies . Told that his claim is nonsense, the man calls a friend who says he too dreamt of the deity the day before yesterday . Two dreams are then treated as proof that the deity exists . That is verifiability and inductive reasoning---and, in Popper's terms, pseudoscience . Religion goes with the verifiability approach .
\end{itemize}
\end{ExampleBox}

\paragraph{The dictated critique of the classical scientific method}
\begin{itemize}
    \item Scientific method itself evolved; it was not developed overnight, but out of multiple experiments .
    \item Karl Popper argued: the problem with the classical scientific method is that it runs on the verifiability approach and inductive reasoning---which cannot be falsified .
    \item Rather, scientific method should follow deductive reasoning and the falsifiability approach .
    \item Pseudo means false: whosoever follows the verifiability-plus-induction route is doing false science .
\end{itemize}

\begin{QuickRevision}
\begin{itemize}
    \item Karl Popper---demarcation principle: separates science from pseudoscience .
    \item Science = falsifiability + deductive reasoning; pseudoscience = verifiability + inductive reasoning .
    \item Classical scientific method is criticised for resting on verification and induction .
    \item Kekulé's benzene ring: a dream became science because it was rigorously tested---a deductive process .
\end{itemize}
\end{QuickRevision}

\subsection{Science and Pseudoscience: Worked Examples}

\vspace{6pt}
\noindent
\begin{tabularx}{\textwidth}{p{3.5cm} p{4cm} Y}
\toprule
\rowcolor{AccentDark}
\thead{Case} & \thead{Verdict given in class} & \thead{Reason} \\
\midrule
\textbf{Astronomy vs. Astrology} & Astronomy is science; astrology is pseudoscience & Only one of the two admits testing and falsification  \\
\textbf{Hypnosis} & Pseudoscience & Claims about unconscious content cannot be tested  \\
\textbf{Crop circles} & Pseudoscience---later revised to ``neither science nor pseudoscience'' & The claim that a spaceship descended and produced the circles can never be tested; the belief exists but cannot be assessed at all  \\
\textbf{Einstein's relativity} & Science & Testable---and therefore falsifiable  \\
\textbf{Newton's gravitation} & Science & Testable; and it was in fact falsified by Einstein  \\
\bottomrule
\end{tabularx}
\vspace{6pt}

\begin{CriticalPoint}
\begin{itemize}
    \item A clarification insisted upon: if you are not able to test something, that does not mean it is nonsense .
    \item There are many theories in science that can never be tested . \textbf{Imre Lakatos}, a critic of Karl Popper, holds that many untestable theories are still part of science .
    \item Quantum physics---we cannot prove what happens there . String theory---we cannot prove that strings exist, yet the community accepts it . Time travel likewise cannot be tested .
    \item Going strictly by Popper's criterion, string theory would have to be called pseudoscience because it cannot be tested .
    \item On yoga, the discussion was deliberately postponed in this session; students were told to keep the Popperian criterion of science versus pseudoscience in mind and they would come very close to the answer without reading any book or journal .
\end{itemize}
\end{CriticalPoint}

\begin{QuickRevision}
\begin{itemize}
    \item Astronomy = science; astrology, hypnosis = pseudoscience .
    \item Crop circles: untestable, and later described as neither science nor pseudoscience .
    \item Lakatos criticises Popper: quantum physics, string theory and time travel are untestable yet remain part of science .
    \item Untestable does not automatically mean nonsense .
\end{itemize}
\end{QuickRevision}

\subsection{Critique of Scientific Method II --- Paul Feyerabend}
The second critic taken up was \textbf{Paul Feyerabend}, who held that we celebrate the scientific method far more than it deserves .

\begin{DataFact}
Approximately 60\% of scientific discoveries are chance experiments---nobody intentionally discovered those things, and nobody had planned them . They were discovered by chance, and only afterwards claimed as achievement .
\end{DataFact}

\begin{ExampleBox}
\textbf{Newton under the apple tree:} At the time the law of gravitation was discovered, there was an epidemic and the university could not be attended---study was from home . One day, bored, he went out and sat under the apple orchard, and an apple fell by accident . Did he follow the scientific method? No . He never set out from home planning to develop a law of gravitation . It was a chance experiment .
\end{ExampleBox}

\begin{KeyConcept}
Such chance discoveries are called \textbf{serendipity}---literally, a happy accident: discovering something beautiful by chance .
\end{KeyConcept}

\begin{ExampleBox}
\begin{itemize}
    \item \textbf{Alfred Nobel:} In whose name the Nobel Prize is given, had never planned to develop gelatin . He cut his finger with a knife and applied a solution to stop the bleeding---the idea behind what we today call a bandage . He then thought that just as a bandage heals a wound and holds two ends together, there could be a substance that holds together and stabilises compounds such as nitrocellulose, making them inactive for the time being until diffused . The actual problem he was working on was why two chemical compounds were not reacting with each other, why they repelled one another . Nothing was planned or systematic; the solution was left overnight and the result was discovered .
    \item \textbf{Louis Pasteur:} Likewise developed a vaccine---the vaccine for chicken pox (or rabies/anthrax), as stated in class---by a chance circumstance .
    \item \textbf{Srinivasa Ramanujan:} Developed a great many mathematical equations, most of which, he said, came to him in a dream, told to him by a goddess; and by chance they turned out to be true . He followed no scientific method---no theorising, observing or data collection .
    \item The psychologist \textbf{Kevin Dunbar} was also named in this connection .
\end{itemize}
\end{ExampleBox}

\paragraph{The dictated propositions of Feyerabend}
\begin{itemize}
    \item Feyerabend argues: science does not go by scientific method . Rather, it goes by the \textbf{``anything goes''} principle . He is clearly mocking and ridiculing science: science does not have any method .
    \item He also believes there is no one scientific method that everybody follows---which means science is characterised by \textbf{methodological pluralism} .
    \item If multiple methods are available, and if more often than not discoveries come without following any method at all, then what is science? How is it different from religion, which also goes by an ``anything goes'' principle? 
    \item Feyerabend argues that to regard science over and above religion is nothing but \textbf{epistemological anarchism} .
\end{itemize}

\begin{CriticalPoint}
A study note given in class: there is no need to go and read up every scholar in detail . Many scholars will be covered, and there will not be time to research each one individually .
\end{CriticalPoint}

\begin{QuickRevision}
\begin{itemize}
    \item Paul Feyerabend---second critique of scientific method .
    \item About 60\% of scientific discoveries are chance experiments---serendipity (happy accident): Newton, Nobel, Pasteur, Ramanujan .
    \item Science does not go by scientific method; it goes by the ``anything goes'' principle .
    \item Methodological pluralism---no single method that all follow .
    \item Ranking science above religion = epistemological anarchism .
\end{itemize}
\end{QuickRevision}

\subsection{Doubts Discussed in Session 2}

\paragraph{1. If Nobel thought about a solution, how can it be chance?}
It was still a chance experiment . He never had the ambition to develop gelatin; the problem he was actually working on was why two compounds---nitrocellulose and another solution---were not reacting and were repelling each other . The bandage analogy occurred to him; the solution was left overnight and the result appeared . Nothing was planned; nothing proceeded in a systematic fashion of the kind we study---review the literature, collect data, classify, then take it to the laboratory . It went very randomly .

\paragraph{2. Chanting in COVID treatment --- is it pseudoscience?}
The question referred to a claim that the effect of chanting on coronavirus patients would be tested . How would you test the effect of chanting?  How would you test the efficacy of a preacher's recitation to someone who has contracted the virus, merely because two days after the sessions the person feels better?  How would you prove that it was only because of the chanting?  These are circumstantial evidences, and science rejects circumstantial evidence . Science needs empirical, testable, objective evidence .

\paragraph{3. Please explain epistemological anarchism}
When you regard science as supreme as a body of knowledge and religion as nonsense, that is epistemological anarchism---because, more often than not, in around 60 to 70\% of cases, science is also based on belief, and most of its experiments happen by chance .

\paragraph{4. Does gravitational force exist? Newton versus Einstein}
Science says yes---but scientific knowledge is always provisional and must be testable . Newton said gravitation exists; Einstein rejected it, and we accept the rejection . There is now a conflict in the scientific community over this, which leads into Thomas Kuhn .
\begin{itemize}
    \item \textit{Newton's position:} The apple fell downward, not upward, because the earth attracted it; similarly all bodies in the universe attract each other, and this explains gravitation .
    \item \textit{Einstein's question to Newton:} If all bodies attract each other, why does the planet Mercury have so different an orbit---as if it did not belong to our solar system?  All other planets show orbits in symmetry; Mercury is deviant . Newton cannot explain this .
\end{itemize}

\begin{ExampleBox}
\textbf{Einstein's answer:} We have a four-dimensional space-time fabric . Place the sun on that fabric and its mass creates a depression---as any object placed on a sheet of paper creates a depression . Because the sun is the heaviest mass, the depression it creates is very heavy, and the ripples from that depression are what displace Mercury's orbit from the simple path expected of it . Gravitation as a force does not exist . What exists is the warping of space and time; the object warps the space-time fabric, and we call the effect a force as though it were external to them .
\end{ExampleBox}

\begin{ComparisonBox}
\begin{itemize}
    \item Newton explained everything in relation to the larger universe: all planets are governed by certain laws; every object is a part of the universe governed by external forces such as the laws of physics .
    \item Einstein showed that each object also has a consequence for the larger universe---the mass of the sun creates the depression that wobbles Mercury's orbit .
    \item The parallel drawn in class: whether your family governs you, or whether you can also change the structure of the family .
\end{itemize}
\end{ComparisonBox}

If you say gravitation is a force, some people will accept you---but after 200 years people will not . Imagine being asked in the Renaissance period whether the earth is at the centre and answering yes; you never knew that after 200 years you would be proven wrong . Similarly, the time of Einstein has come; Newtonian physics is being said goodbye to . That shows that science does not give us objective results: the results and findings of science are subject to the scientific community and their acceptance . Today many young scientists are in love with Einstein while the older generation is still very much with Newton---a conflict to be taken up with Thomas Kuhn .

\paragraph{5. Does a discovery made by accident follow inductive reasoning?}
Not even inductive reasoning---forget the concept of reasoning altogether .

\begin{ExampleBox}
You are sitting at a café having a cold coffee and garlic bread, and out of nowhere a school friend you have not seen for ten or twelve years walks by . That was by chance; it was not planned . Is there any reasoning involved in it---inductive or deductive? No, nothing . Everything is by chance . This is what happens with scientists: many discoveries happen by chance---these are chance experiments (serendipity) .
\end{ExampleBox}

\paragraph{6. Since scientific knowledge is provisional, is it reliable?}
If something is provisional, it is not permanent; if it is not permanent, it cannot be called reliable . It can, however, be tested---or it may not even be tested .
\begin{itemize}
    \item When would knowledge be called absolute rather than provisional? In earlier times people said God exists and must be obeyed; there was no challenge and no question, and the knowledge was believed to be absolute .
    \item Today nothing can be asserted with absolutism . ``The earth exists'' is not provisional; but ``an Earth-like planet also exists in our universe'' is provisional---it is yet to be tested, and therefore not reliable .
    \item It also depends on who is conducting the research and what findings they come up with---including who funds them and what interests those funders have . A Marxist reading would say it is the material economic conditions that explain why a particular scientific exploration happened at all, and whose interests it secures . In the social world nothing is simple .
\end{itemize}

\paragraph{7. Newton's theory was rejected by Einstein --- then how is it scientific?}
Scientific does not mean absolute truth . Scientific only means that it can be tested . We tested Newton's theory; it came out to be false---Einstein did it; it was not really a force but the warping of space and time . Would you still call Newton scientific? Yes---precisely because we could test it . Nothing is scientific because it is absolute or right; it is scientific only if it is testable .

\begin{QuickRevision}
\begin{itemize}
    \item Chance discovery involves no reasoning at all---neither inductive nor deductive .
    \item Chanting claims rest on circumstantial evidence, which science rejects .
    \item Newton $\rightarrow$ Einstein: Mercury's orbit unexplained; gravity is not a force but the warping of space-time .
    \item Science's findings depend on the acceptance of the scientific community---the bridge to Thomas Kuhn .
    \item Provisional knowledge is not reliable---but it remains scientific as long as it is testable .
\end{itemize}
\end{QuickRevision}

\sectiondivider

\section{Provisional Knowledge, Kuhn, Feminist Critique \& Theoretical Strands}

\subsection{Testing Everyday Statements: Light, Sound and Photons}
The session began where the previous one ended---with the critique of scientific method---by putting a series of ordinary statements to the test of testability .

\paragraph{``Light travels faster than sound''}
The statement is scientific because it can be tested .

\begin{ExampleBox}
It is raining; suddenly a huge thunderbolt is seen, and only afterwards the sound of the burst is heard (the gap was exaggerated in class to a minute; in fact it is five or six seconds) . The bolt itself was the result of the collision between clouds, but the light was seen first . That clearly shows light travels faster than sound---something you can experiment with yourself, which is why it is a scientific statement, and it has been proven right .
\end{ExampleBox}

\paragraph{``Nothing travels faster than light''}
Asked 50 years ago, everyone would have answered yes . Asked today, the class was confused .

\begin{ExampleBox}
A religious teacher, told this statement, replied that it is wrong: God is everywhere . If God is everywhere, then God is travelling faster than light---how else would you explain presence at two different places at the same point of time?  On the one hand we say nothing travels faster than light; on the other we say God is present everywhere---and we accept both as valid .
\end{ExampleBox}

Today the statement itself is doubted: quantum physics tells us sub-atomic particles travel faster than light . But has it been proven? Not yet . If such a question were asked from a textbook it would be a chaotic situation; you could not reach an absolute answer, which is why it would not be asked . And if such a statement is printed, after four or five years the textbooks will be revised---new knowledge given, old knowledge discarded .

\paragraph{``Light consists of particles''}

\vspace{6pt}
\noindent
\begin{tabularx}{\textwidth}{p{3.5cm} p{4.5cm} Y}
\toprule
\rowcolor{AccentDark}
\thead{Position} & \thead{Who / When} & \thead{Status} \\
\midrule
\textbf{Light consists of particles} & Newton, in his book \textit{Optics}, 1704 & Old knowledge; rejected and removed from textbooks  \\
\textbf{Light consists of waves} & Thomas Young---proved through the double slit experiment & Replaced the particle view  \\
\textbf{Light consists of photons} & Quantum physics---the latest theory & Neither particles, nor waves, nor both  \\
\bottomrule
\end{tabularx}
\vspace{6pt}

Which of them is right---waves, particles, both, photons, all of the above, or none? No answer can be given . Tomorrow someone will cancel this thesis too and a new knowledge will develop .

\begin{CriticalPoint}
\begin{itemize}
    \item The point being pressed: science is not objective or absolute, and it does not have a monopoly over truth, though it claims one .
    \item If the answers given by scientists keep changing, are we not studying myth in the name of science? The question was put to the class in exactly that form .
    \item The person who studied the earlier position finds that training rendered obsolete (the teacher noted having completed a technical degree and lost belief in science, feeling that only wrong information had been learnt) .
\end{itemize}
\end{CriticalPoint}

\begin{QuickRevision}
\begin{itemize}
    \item Light travels faster than sound---testable, hence scientific, and proven .
    \item ``Nothing travels faster than light'' is now contested by quantum physics (unproven) .
    \item Particles (Newton, \textit{Optics}, 1704) $\rightarrow$ waves (Thomas Young, double slit experiment) $\rightarrow$ photons (quantum physics) .
    \item Science has no monopoly over truth; textbooks are revised as knowledge changes .
\end{itemize}
\end{QuickRevision}

\subsection{Falling Bodies: Aristotle to Einstein}
A mobile phone and a pen were held up: dropped together, which falls first?  The history of that one question was then traced .

\vspace{6pt}
\noindent
\begin{tabularx}{\textwidth}{p{3cm} p{4cm} Y}
\toprule
\rowcolor{AccentDark}
\thead{Scholar} & \thead{Claim} & \thead{Explanation offered} \\
\midrule
\textbf{Aristotle} & The heavier object falls first & The heavier object has innate, natural, internal properties that make it fall before the lighter one  \\
\textbf{Galileo Galilei} & Both fall together, irrespective of mass & Performed the experiment from the Tower of Pisa; both attain a specific constant acceleration  \\
\textbf{Isaac Newton} & Both fall together & Explained by the force of gravity, which is external to the masses and exerts a pull on them  \\
\textbf{Albert Einstein} & Gravity as a force does not exist & Only the warping of space and time  \\
\bottomrule
\end{tabularx}
\vspace{6pt}

At each stage the earlier authority was rejected: Aristotle was rejected by Galileo, who did not simply listen to Aristotle but performed an experiment; Galileo's acceleration was recast by Newton as force; Newton's force was dissolved by Einstein . Is Einstein correct? Today you think so . Tomorrow he too may be rejected . Will you ever be in a position to say who is right and who is wrong? 

\begin{QuickRevision}
\begin{itemize}
    \item Aristotle (innate properties) $\rightarrow$ Galileo (Tower of Pisa; constant acceleration) $\rightarrow$ Newton (force of gravity) $\rightarrow$ Einstein (warping of space-time) .
    \item Every dominant answer was overturned by the next---including, possibly, the present one .
\end{itemize}
\end{QuickRevision}

\subsection{Gravitational Waves and the Tests That Confirmed Einstein}
Imagine the universe as a sheet of paper . Place the sun on it and a bend is created---as any mass placed on a sheet creates a bend .
\begin{itemize}
    \item That bend creates ripples, exactly as a stone thrown into a lake creates ripples . All the other objects---planets, stars, asteroids---feel the consequence of those ripples, and their positions are displaced accordingly .
    \item So it is not a downward motion or a pull maintaining the balance of the universe: it is a wave---a \textbf{gravitational wave} . This is what Einstein told us before he died .
\end{itemize}

\begin{DataFact}
\begin{itemize}
    \item \textbf{The LIGO experiment:} An L-shaped tunnel, a very heavy apparatus, built specifically to verify what Einstein had said . The year was recalled as around 2016 (possibly 2016, 2017 or 2018---stated as uncertain in class) . Einstein was right; the waves were found .
    \item \textbf{The 1919 solar eclipse experiment:} If the space-time fabric is bent by mass and energy, then light thrown across the universe must also bend . That light really bends was proven in 1919 .
\end{itemize}
\end{DataFact}

\begin{KeyConcept}
\begin{itemize}
    \item First conclusion of the story: a theory is scientific only when it is testable .
    \item Second conclusion: no scientific knowledge is absolute---these are only provisional, temporary statements .
    \item And who taught us this lesson? \textbf{Karl Popper} .
\end{itemize}
\end{KeyConcept}

\begin{QuickRevision}
\begin{itemize}
    \item Sun on the space-time sheet $\rightarrow$ bend $\rightarrow$ ripples $\rightarrow$ gravitational waves .
    \item Confirmed by the LIGO experiment (around 2016) and, for the bending of light, the 1919 solar eclipse experiment .
    \item Scientific = testable; scientific knowledge = provisional . Lesson from Karl Popper .
\end{itemize}
\end{QuickRevision}

\subsection{Is Marxism Scientific or Pseudoscience?}
A statement was put to the class: the whole economic system will collapse one day due to inherent contradictions in the economy; societies will be ruled by workers; the industrialists will be overthrown by the workers, who have been exploited enough, and the working class will decide the policy of the state . Can you test this statement? No . Therefore, judged by the criterion in hand, Marxism is pseudoscience .

\begin{CriticalPoint}
\begin{itemize}
    \item But if you told Marx that his theory is pseudoscientific, he would say the same of yours; and if you sat with him for ten hours you would come to believe his theory is scientific .
    \item When you called his theory pseudoscience you were looking at it through the lens of Karl Popper . Remove that lens, put on the lens of Marxist scholars, and the theory becomes scientific .
    \item So is there any objective reality in this world? Perspectives only .
\end{itemize}
\end{CriticalPoint}

\begin{QuickRevision}
\begin{itemize}
    \item The Marxist prediction of collapse cannot be tested $\rightarrow$ pseudoscience by Popper's criterion .
    \item The verdict depends on the lens used---there are only perspectives .
\end{itemize}
\end{QuickRevision}

\subsection{Thomas Kuhn --- Science as a Social Process}
\textbf{Thomas Kuhn} was introduced as the biggest figure among these thinkers---a sociologist of science, and also a philosopher and historian of science . Note that he is a philosopher of science as a body of knowledge, not merely of natural science .
\begin{enumerate}
    \item \textbf{First proposition (dictated):} \textit{``Science is a social process.''}  It was noted in class that this statement itself can appear in the examination in coming years .
    \item \textbf{Second proposition (dictated):} Any theory is valid only so far, until it is superseded by another one . That is, the validity of theory is provisional . There is no absolute theory: after some point this theory will be rejected by another, and people will then treat the new one as more scientific, or as the absolute truth---and this keeps happening .
    \item \textbf{Third proposition (dictated):} A scientist operates within a \textbf{dominant paradigm}---a scientist meaning both natural scientists and social scientists .
\end{enumerate}

\begin{KeyConcept}
\begin{itemize}
    \item \textbf{Dominant paradigm} = any set of theory or belief which is accepted as absolutely true, and which has rejected the theory that existed before it .
    \item Until Einstein rejected Newton, the dominant paradigm was Newtonian physics .
    \item The technical vocabulary---paradigm, dominant paradigm---was deliberately kept simple at first so that the idea is understood before the term .
\end{itemize}
\end{KeyConcept}

\begin{ExampleBox}
Before Galileo rejected Aristotelian physics, Aristotelian physics was the dominant paradigm; in those times people read that heavier objects fall first . After Aristotle came Galileo; after Galileo, Newton---and the dominant paradigm of the world was gravitational force . Today the dominant paradigm is Einstein . But this transition is not smooth . You will always face chaos . If after this class you go out and tell someone that gravity as a force does not exist, he will laugh at you . A dominant paradigm is valid for a very limited period; it has an expiry date . It expires when a new perspective develops which counters it .
\end{ExampleBox}

\begin{QuickRevision}
\begin{itemize}
    \item Thomas Kuhn---sociologist, philosopher and historian of science .
    \item Science is a social process; any theory is valid only until superseded; a scientist operates within a dominant paradigm .
    \item Dominant paradigm = a set of theory or belief accepted as absolutely true, having rejected its predecessor; it expires when countered .
\end{itemize}
\end{QuickRevision}

\subsection{The Structure of Scientific Revolutions}
Kuhn's book is \textit{The Structure of Scientific Revolutions} . The revolution he speaks of is an upheaval in the world of science---a belief system changing wholly, a paradigm collapsing . Newton is gone, and with him the theory that gravity is a force; Einstein has arrived . That is revolution .

\begin{ExampleBox}
When the geocentric theory was first challenged by the heliocentric model: the geocentric view had been given by Greek astronomers---Ptolemy and Aristotle---and the church accepted what they said . These were astronomers, not astrologers, and astronomy we regard as scientific . It was rejected when Copernicus came . But until Copernicus came, was it not the dominant paradigm? Yes it was---and in Kuhn's language, for that period it was \textbf{normal science} . Just as today we accept Einstein's theory as correct and Einsteinian physics is normal science, at that time geocentric theory was normal science .
\end{ExampleBox}

\paragraph{The Kuhnian Cycle}
$$\text{Normal Science} \longrightarrow \text{Crisis} \longrightarrow \text{Revolutionary Science} \longrightarrow \text{New Normal Science}$$
\begin{enumerate}
    \item \textbf{Normal science:} The accepted body of theory of the time .
    \item \textbf{Crisis:} People begin criticising that so-called normal science .
    \item \textbf{Revolutionary science:} Emerges out of the crisis .
    \item The revolutionary science becomes the new normal science; it will again be rejected, crisis will follow, and a new revolutionary science will develop .
\end{enumerate}

\begin{CriticalPoint}
Why is the transition violent? If you have developed a theory over ten years and someone appears from nowhere to say you are wrong, you will not accept the criticism easily . This is why the church declared Galileo Galilei and Copernicus heretics . None of these people accepted the challenge easily .
\end{CriticalPoint}

\begin{ComparisonBox}
A parallel from German philosophy: Hegel says the world is nothing but the product of \textbf{dialectics between ideas}---one idea exists, another comes, the first goes away, the new one dominates, and next year yet another replaces it . This is how the world progresses and how we come to know what is true . This is known as the law of dialectics .
\end{ComparisonBox}

\begin{PYQLink}
All this is useful beyond sociology---particularly in the essay paper, where one essay must compulsorily be written on a philosophical topic, and topics such as science versus religion recur . Command over the history of science or the philosophy of science gives you the whole background for such an essay .
\end{PYQLink}

\begin{QuickRevision}
\begin{itemize}
    \item Book: \textit{The Structure of Scientific Revolutions} .
    \item Cycle: normal science $\rightarrow$ crisis $\rightarrow$ revolutionary science $\rightarrow$ new normal science .
    \item Geocentrism (Ptolemy, Aristotle, backed by the church) was normal science until Copernicus; Galileo and Copernicus were declared heretics .
    \item Parallel: Hegel's law of dialectics .
\end{itemize}
\end{QuickRevision}

\subsection{Who Decides the Dominant Paradigm?}
The question was put directly: who decides that Einstein is right and Newton wrong? Not you and not I . The \textbf{scientific community} decides .
\begin{itemize}
    \item What is the scientific community? Institutions such as IEEE, NASA, ISRO, DRDO---the bodies which decide whether Newton or Einstein is right .
    \item Only after that decision is a communication sent to the textbook authorities---include this chapter in class eight from next year onwards, the content has been revised---and the new position enters the textbooks of school students, who read it as settled fact .
\end{itemize}

\begin{CriticalPoint}
How is the decision actually taken? In class it was described satirically: they sit together in a closed, air-conditioned chamber, take many cups of tea and plates of snacks, deliberate on whether Einstein seems scientific, and finally decide to go ahead with Einstein . From the next day onwards Einstein is the father of physics and the dominant paradigm . After twenty years, when asked whether gravity exists, students will answer that gravity as a force does not exist, it is the warping of space---and by then a new scientist may have rejected warping too .
\end{CriticalPoint}

\begin{KeyConcept}
\begin{itemize}
    \item \textbf{Dictated:} \textit{Unless the scientific community decides and agrees upon their decision, the theory is not considered scientific.} 
    \item Is science not a social process, then? It is---and not only because a community sits on the fate of science, but also because science is very temporal and spatial: what was valid a hundred years back is not valid today .
    \item Just as the human body matures and grows, science also grows . Every day scientists refute the theses in which they had reposed faith for years, and come forward with a new one on which all of them agree; that then becomes part of your syllabus, while we are taught that science is highly objective and absolute . There is no absolute truth; all knowledge is provisional .
\end{itemize}
\end{KeyConcept}

\begin{QuickRevision}
\begin{itemize}
    \item The scientific community (IEEE, NASA, ISRO, DRDO) decides what counts as scientific; textbooks follow .
    \item Dictated: unless the scientific community decides and agrees, a theory is not considered scientific .
    \item Science is a social process because it is community-sanctioned and temporally and spatially limited .
\end{itemize}
\end{QuickRevision}

\subsection{Feminist Critique of Scientific Method}
The last critique of scientific method taken up in this session was the feminist one . The scholars named were \textbf{Donna Haraway}, \textbf{Sandra Harding} and \textbf{Evelyn Fox Keller} .

\begin{KeyConcept}
\begin{itemize}
    \item \textbf{The feminist claim:} Science, or scientific method, is sexist and patriarchal---it reaffirms or reproduces sexism in society .
    \item Two disciplines are especially in view because both call themselves scientific: \textbf{biology} and \textbf{psychology} .
\end{itemize}
\end{KeyConcept}

\paragraph{Biology I --- Active sperm, passive egg}
\begin{itemize}
    \item Biology speaks of the ``active sperm''---the term is still in use . Why should the sperm be called active? And when the egg is called passive, the egg is displayed as sitting there, waiting to be reached---projecting a docile woman .
    \item This is a sexist belief among scientists which makes them write such terms for sperm and egg; it is the reproduction of the patriarchal mindset scientists carry .
    \item Some may say the sperm is active because it has to travel . That theory too was rejected---as stated in class, the sperm does not travel .
\end{itemize}

\paragraph{Biology II --- Charles Darwin}
\begin{itemize}
    \item Darwin wrote \textit{The Descent of Man} . Darwin is regarded as a very important scholar; without him there would have been no subject like biology .
    \item Darwin held that in the course of evolution man naturally evolves superior to women---this is how societies evolve naturally . Everything in the claim rests on that word .
    \item After reading ``natural, natural selection, natural'' repeatedly, you begin to reproduce the same knowledge . It becomes part of your common sense, and you repeat it to your friends: have you not read Darwin---it is natural . Fox Keller and Donna Haraway reply that nothing is natural .
    \item Darwin also said that if two lists were made, of men and of women, in artistry, poetry, drama and theatrics, the lists would bear no comparison, and man would emerge superior---in sport, literary work, dance and music alike, men dominate .
\end{itemize}

\begin{ExampleBox}
\textbf{Antoinette Brown Blackwell}---described in class as the first feminist critic of Darwin, officially---wrote a letter to Darwin about his theories, telling him they were nonsense and highly patriarchal . Darwin replied: \textit{``Dear Mr. Antoinette.''}  What does that show? That in his mind a woman could not have sent such a letter, could not have criticised his scientific theories---it would have to be a man . That is the so-called ``natural'' at work . In her next letter she pointed out that she is not Mr but Mrs---and therefore his ``natural'' thesis is rejected .
\end{ExampleBox}

\paragraph{Psychology --- Sigmund Freud and Carl Jung}
\begin{itemize}
    \item Sigmund Freud gave the theory of the \textbf{Oedipus complex}: when a male child is born, he begins to feel sexual attraction towards his mother . Is this scientific by Popper's criterion? Can it be tested that as a two-year-old male child you felt sexual attraction towards your mother? No .
\end{itemize}

\begin{CriticalPoint}
\begin{itemize}
    \item That is why, to Karl Popper, Freud and Marx were the biggest enemies of science---and both called their theories scientific: Marx called his \textit{scientific socialism}, Freud's work is generally regarded as \textit{scientific psychoanalysis} . Popper rejected both .
    \item And where is the father while this is happening? Freud says the father is killed, avenged or hated by the male child .
    \item \textbf{The feminist reading:} In this story the mother is absent . She is projected as someone simply waiting for the male child to feel attraction---as though she will not scold him, will not socialise him, will not train him, and sits idle .
    \item Women are projected as objects without a voice of their own . Reading Freud you always read an antagonistic relation between the male child and the father---a fight over the mother, as if the mother is an object to be claimed by one or the other . This normalises the belief in society that a woman is something to be caught or fought for . That is the feminist critique of scientific method .
    \item \textbf{Carl Jung} then gave the \textbf{Electra complex}---the exact opposite of the Oedipus complex: the female child feeling sexual attraction towards the father . Jung thereby criticised Freud by reversing the whole theory, which could be equally valid .
\end{itemize}
\end{CriticalPoint}

\paragraph{Numbers, socialisation and detached observation}
\begin{itemize}
    \item Look at the IITs, DRDO, ISRO or NASA: how many women would you find? The norm of universalism says social location does not matter, yet in the real world women are not highly active in the sciences . Science does discriminate .
\end{itemize}

\begin{ExampleBox}
A recollection offered in class: through twelve years of school all the teachers were women, so one was socialised into believing that teachers in society are female . On joining college, seeing a male teacher was itself disorienting .
\end{ExampleBox}

\begin{CriticalPoint}
\begin{itemize}
    \item The last feminist critique concerns the values enshrined in the scientific method: objectivity, empiricism, experimentation, rationality .
    \item When you give too much importance to rationality, objectivity and empiricism, you lack empathy---you lack emotions and the human values that are natural .
    \item \textbf{Personality test example:} The personality test taken after the mains is generally regarded as a scientific method of evaluating personality, and is believed to yield the best assessment . Feminists would say the method being employed is highly patriarchal: you are asked to read ethics in order to develop empathy, and then are not expected to manifest empathy or be expressive; you are wanted to be serious rather than smiling . In such an encounter you feel exactly what women feel in the ordinary power relation between men and women .
    \item So the scientific method lacks human values, and what it produces is a \textbf{detached observation}---the person interested in knowing another person stays highly detached from him in order to maintain so-called objectivity . And detached observation is not observation .
\end{itemize}
\end{CriticalPoint}

\begin{QuickRevision}
\begin{itemize}
    \item Feminist scholars: Donna Haraway, Sandra Harding, Evelyn Fox Keller---science is sexist and patriarchal and reproduces sexism .
    \item Biology: active sperm / passive egg; Darwin's ``natural'' superiority of man---answered by Antoinette Brown Blackwell (``Dear Mr. Antoinette'') .
    \item Psychology: Freud's Oedipus complex (untestable; Popper rejected Freud and Marx); the mother as a voiceless object; Jung's Electra complex .
    \item Institutional under-representation of women; and the method's exclusion of empathy---detached observation is not observation .
\end{itemize}
\end{QuickRevision}

\subsection{Yoga: Science or Pseudoscience?}
The question deferred earlier was answered when a student pressed it . To answer it there is no choice but to take the help of Karl Popper .
\begin{itemize}
    \item Yoga is not only about the exercise we do; yoga has multiple philosophies, and there is a Ministry of AYUSH covering Ayurveda, Yoga, Unani, Siddha, Sowa-Rigpa and Homoeopathy .
\end{itemize}

\begin{ExampleBox}
During the lockdown, people collectively banged plates and lit candles at an appointed hour . The stated rationale was that this would generate a collective consciousness and help us move beyond that phase---a yogic practice of collective consciousness . Durkheim's collective consciousness, to be taught later, is the concept in the background . Compare the temple: you come out feeling relaxed and peaceful---not because of the building but because of the collective consciousness of that gathering and the act of closing your eyes together . An article was referred to discussing Popper's idea and how we failed as a scientific age, since in a scientific age we believed such practices could carry us over a chaotic phase . This was pseudoscience .
\end{ExampleBox}

\begin{CriticalPoint}
Even taking yoga purely as exercise: can you ever prove that you were healed only because of the exercise, with no other external factor at work? You would not be able to prove it . Therefore, to Karl Popper, yoga is not science---and this is why a scientific body, the Indian Medical Association, says the same: it is pseudoscience .
\end{CriticalPoint}

\begin{QuickRevision}
\begin{itemize}
    \item Judged by Popper: yoga is pseudoscience---the healing claim cannot be isolated from external factors and cannot be falsified .
    \item Yoga is more than exercise; it carries multiple philosophies (Ministry of AYUSH) .
    \item Lockdown collective rituals framed as collective consciousness---Durkheim in the background .
\end{itemize}
\end{QuickRevision}

\subsection{Major Theoretical Strands: The Macro or Structural Approach}
Sub-unit 2B is \textbf{major theoretical strands of research methodology}---that is, the major theoretical perspectives of research methodology (which overlaps with sub-unit 2E) . The first strand is the \textbf{macro or structural approach} .
\begin{itemize}
    \item If I say that family is a unit that helps society maintain itself, that it socialises the young and perpetuates the norms and values required for socialisation---then I am speaking as a structural functionalist .
    \item If I say that religion is a unit which helps society get over a crisis situation---as when, during the lockdown, people became far more religious, mythological serials were aired, and forwarded messages and statuses turned to matters of faith---that is the same approach .
    \item In the macro or structural approach you always talk about the function of each and every unit in the maintenance of society .
\end{itemize}

\paragraph{The three dictated points}
\begin{enumerate}
    \item Structural approach is an approach which \textbf{regards structure as superior to individuals} .
    \item It also projects \textbf{individuals as passive robots}, because no agency is given to them; all the structures in society govern society and the individuals .
    \item This approach \textbf{fails to look at how individuals attach meaning} to these institutions .
\end{enumerate}

\begin{ExampleBox}
Asked about the function of marriage, the standard answer given is: maintenance and continuity of society . But consider someone marrying in order to obtain a passport . In ancient Indian empires the practice of marriage was used to make political alliances with other rulers---not to continue society . Suppose you become an administrative officer: a legislator may seek to marry his daughter or son to you . How is he looking at marriage? As a chance for his child to improve their standard of living, and as a chance to forge an alliance with the bureaucracy . And you may accept, seeing it as a chance to enter politics . How we use and look at the institution of marriage is what this approach does not talk about; it talks only of the function of marriage, and of religion, to the system .
\end{ExampleBox}

\begin{QuickRevision}
\begin{itemize}
    \item Macro/structural approach = the function of each unit in the maintenance of society .
    \item Three points: structure superior to individuals; individuals as passive robots; fails to see how individuals attach meaning to institutions .
    \item Marriage, religion and law are read only through their function to the system .
\end{itemize}
\end{QuickRevision}

\subsection{Positivism}
The most dominant of the structural approaches is \textbf{positivism}, or positivist methodology . Being a structural approach, the three points above apply to it as well; what follows is what is special to it :
\begin{enumerate}
    \item Positivism assumes society is governed by laws . The assumption came from the influence of Enlightenment philosophers . Auguste Comte was in fact very critical of the Enlightenment scholars, but one thing that inspired him was their talk of reasoning---from which he went on to speak of laws operating in the social world as well .
    \item These laws govern how individuals, groups or societies think in each era .
    \item \textbf{Law of three stages:} In each era individuals, groups or societies think according to the law---the theological stage (religious beliefs), the metaphysical stage (abstract beliefs), and the positive stage (scientific beliefs) .
    \item They also regard each unit as contributing to the functioning of the system .
\end{enumerate}

\begin{QuickRevision}
\begin{itemize}
    \item Positivism = the dominant structural approach .
    \item Society is governed by laws; those laws govern how each era thinks; law of three stages---theological, metaphysical, positive; each unit contributes to the system .
    \item Comte was critical of the Enlightenment yet inspired by its emphasis on reasoning .
\end{itemize}
\end{QuickRevision}

\subsection{Structural Functionalism}
Three points were dictated for structural functionalism :
\begin{enumerate}
    \item Structural functionalism, being a positivist theory, also believes society is governed by laws .
    \item \textbf{Which law?} Herbert Spencer is generally regarded as the father of structural functionalism, and through the \textbf{organic analogy}---human bodies and societies are very much similar, and just as the human body evolves, society also evolves---the law operating here is the \textbf{law of evolution} . Society is governed by the law of evolution, which implies that each structure of society also evolves .
    \item This theory largely believes that each social institution or structure contributes to the maintenance of the system, which also implies that each institution plays a \textbf{positive function only}---the positive function being that the maintenance of society is supreme .
\end{enumerate}

\begin{ExampleBox}
\begin{itemize}
    \item \textbf{Evolution of structures:} The family disintegrating from joint to nuclear .
    \item \textbf{Religion:} In earlier times the church disintegrated, and today we have fragments---cults, spiritual figures, and books on elevating the mind or the power of the mind, through which you can become your own guru or your own religion . Religion is no longer found in the church alone: it appears on television at five in the morning and again in the evening . This too is a fragmented form of religion .
\end{itemize}
\end{ExampleBox}

\begin{CriticalPoint}
Do not bring Robert K. Merton into this discussion at this stage . Here we are discussing structural functionalism as the dominant paradigm of one period; Merton is exceptional and is treated separately .
\end{CriticalPoint}

\begin{QuickRevision}
\begin{itemize}
    \item Structural functionalism is a positivist theory---society governed by laws .
    \item The law here is the law of evolution; Spencer as father of the school; organic analogy .
    \item Each institution contributes to the maintenance of the system and therefore plays a positive function only .
\end{itemize}
\end{QuickRevision}

\subsection{Marxism as a Positivist Philosophy}
\begin{itemize}
    \item Marxism is also a positivist philosophy, and being a positivist philosophy it also believes that society is governed by laws .
    \item However, unlike structural functionalists, Marxism believes that institutions function not for the maintenance of society but \textbf{to secure the interest of the capitalist class} .
    \item \textbf{Which law governs society according to Marx?} The \textbf{law of dialectics} .
\end{itemize}

\begin{CriticalPoint}
\begin{itemize}
    \item Do not say dialectical materialism here---dialectical materialism is also right, but the word required at this point is law .
    \item The laws of dialectics were postulated by \textbf{Friedrich Engels}---not Hegel---and Marx applied these laws .
    \item Some argue Marxism is only partially positivist . There are many journals on this debate; it is deliberately left aside so that the discussion does not become endless .
    \item A chart placing all these strands together was made in class to help locate exactly where each perspective belongs .
\end{itemize}
\end{CriticalPoint}

\begin{QuickRevision}
\begin{itemize}
    \item Marxism is positivist: society governed by laws---the law of dialectics (postulated by Engels, applied by Marx) .
    \item Difference from structural functionalism: institutions secure the interest of the capitalist class, not the maintenance of society .
\end{itemize}
\end{QuickRevision}

\subsection{The Micro Approach and Non-Positivism}
\begin{enumerate}
    \item Unlike the macro approach, this approach emphasises the \textbf{role of individuals} . The examples of marrying a legislator's son or daughter belong here---the concern is with you, not with institutions governing society .
    \item However, that does not mean that this approach rejects structures altogether . It is just that they place more importance on the role of individuals .
    \item Therefore, individuals are portrayed not as passive robots but as \textbf{active players} .
\end{enumerate}

\begin{KeyConcept}
\begin{itemize}
    \item The theoretical perspectives dealt with in the micro or action approach are called \textbf{non-positivism}---non-positivist methodologies .
    \item If in positivism the first line was ``society is governed by laws'', then here we write: \textit{society is not governed by laws} . Rather, individuals in society are actively shaping their lives . We are not simply objects governed by the laws of physics .
\end{itemize}
\end{KeyConcept}

\begin{QuickRevision}
\begin{itemize}
    \item Micro approach: role of individuals emphasised, structures not rejected outright .
    \item Individuals = active players, not passive robots .
    \item The perspectives here are called non-positivist methodologies; society is not governed by laws .
\end{itemize}
\end{QuickRevision}

\subsection{Social Action Theory --- Max Weber}
\begin{itemize}
    \item Max Weber gave the theory of social action, and he believed that it is the \textbf{meanings and motives} of the individual actors that should be studied by sociologists .
    \item If every actor has his own meaning and motive behind any action, then individuals are clearly not governed by laws: they have their own active mindset, and behind everything they do there is a specific meaning and motive which only they know .
    \item The theory was given by Weber so as to study the meanings and motives of the individual actors . However, Weber believed that sociology should study not only the cause of social action but also its \textbf{consequences} .
\end{itemize}

\begin{KeyConcept}
\begin{itemize}
    \item \textbf{Cause of social action} = the individual's meanings and motives .
    \item \textbf{Consequence of social action} = structure .
    \item Diagram: $\text{Many individuals with meanings/motives (cause)} \longrightarrow \text{Social Action} \longrightarrow \text{Structure (consequence)}$ .
\end{itemize}
\end{KeyConcept}

\begin{ExampleBox}
\textbf{How did India become a nation state?} At one point of time it was the British Indian Empire; we were not the nation state we are today . The nation state is a structure---a modern structure . It was built up by the social action of the individuals involved in the nationalist movement: their involvement, and the meanings and motives behind that involvement, led to the formation of the nation state . When millions of individuals are mobilised through social action, large-scale structural change occurs .
\end{ExampleBox}

\begin{QuickRevision}
\begin{itemize}
    \item Weber's social action theory: study the meanings and motives of individual actors .
    \item Cause of social action = meanings and motives; consequence of social action = structure .
    \item Illustration: the nation state as the structural consequence of the nationalist movement's mass social action .
\end{itemize}
\end{QuickRevision}

\subsection{Phenomenology --- Alfred Schutz}
Phenomenology developed after the social action approach---remember the sequence: first came Max Weber .
\begin{itemize}
    \item Phenomenology is mainly a philosophical school, developed by \textbf{Edmund Husserl} in philosophy . It was introduced into sociology by \textbf{Alfred Schutz}, who was Husserl's student .
\end{itemize}

\begin{ExampleBox}
\textbf{The red rose:} If I give a red rose to a girl, it means I am proposing to her---this is what we understand, because it is what we have seen around us and in films . The important question is: when and how did we come to know this, and why did we come to believe it---that red means love, that red equals proposal?  We take it for granted . It is part of our common sense . And sociology challenges what is taken for granted . Give a purple rose instead and nobody can say what you are trying to do; we get confused, because we do not know the meaning attached to it .
\end{ExampleBox}

\begin{KeyConcept}
\textbf{Phenomenology is an epistemological school according to which there is no objective reality in the social world.} 
\end{KeyConcept}

\begin{ExampleBox}
\begin{itemize}
    \item \textbf{Bargaining:} In the street vegetable market you are told a price and you bargain, calling it costly . In a shopping mall the same or a higher price is paid without a second's bargaining, cash handed over or card swiped . Is there any objective, rational law that explains this behaviour? These are only socially constructed norms .
    \item \textbf{Classroom as illusion:} The classroom runs because you carry the norm that when the teacher is explaining, students are silent; the moment that norm is disrupted, the classroom as a structure is not possible . If instead of teaching we start playing cricket in it, the same room converts into a cricket ground . It is a classroom only so long as all of us agree that it is .
    \item \textbf{The officer's social order:} When you become an administrative officer and sit in the official car, someone opens the door and the gate for you and greets you . Imagine nobody there: you would have to open the door yourself . You carry the order in your mind that as an officer you have all these privileges---but this is illusion, created so that you keep believing you are an officer . The moment those people are gone, you are not . Along with the bungalow, the lawn, the gardener: this is social order, constructed to make you believe .
\end{itemize}
\end{ExampleBox}

\begin{CriticalPoint}
\begin{itemize}
    \item The postmodern scholar \textbf{Jean Baudrillard} holds that in the age of information you cannot differentiate between content and advertisement---advertisement becomes content .
    \item \textbf{The line to write:} \textit{Social world is an illusion. What we call social order becomes possible only when all of us agree upon it and share it.} 
    \item When the red rose made sense to you, it was because you agree that red means love . The same red on the road means stop---there you do not read the red as a proposal from the traffic policeman . The meaning of anything is not objective; it is highly contextual .
    \item Psychological thrillers such as \textit{Mulholland Drive}, \textit{Shutter Island} and \textit{Inception} fascinate us because they destroy our social order---they shatter the way we think in the real world, and that is why they make us anxious .
\end{itemize}
\end{CriticalPoint}

\begin{QuickRevision}
\begin{itemize}
    \item Phenomenology: philosophical school of Edmund Husserl, brought into sociology by Alfred Schutz (his student); developed after Weber's social action .
    \item Epistemological school: there is no objective reality in the social world .
    \item Social world is an illusion; social order exists only because all of us agree upon and share it .
    \item Meaning is contextual, not objective---red as love, red as stop .
\end{itemize}
\end{QuickRevision}

\sectiondivider

\section{Non-Positivist Methodologies: Phenomenology, Ethnomethodology, Symbolic Interactionism \& Dramaturgy}

\subsection{Phenomenology Restated: The Four Propositions}

\begin{ExampleBox}
\textbf{An ordinary evening:} You return home after work or study . Your parents ask how the preparation is going; you say everything is perfect . Your mother asks what you would like to eat; you say anything . You eat, you say it was nice . You study, call friends, watch videos, and sleep with a target for 5 a.m.  Did you find anything bizarre or abnormal? No . Everything is very normal---it is part of common sense . And yet, according to phenomenologists, none of these things was objective reality . Every conversation, the food on the table, the bed made ready---all of it looks orderly, and all of it is illusion .
\end{ExampleBox}

\paragraph{The four propositions to be written}
\begin{enumerate}
    \item \textbf{Objectivity does not exist} in the social world .
    \item \textbf{Social order is an illusion} .
    \item We tend to \textbf{typify} phenomena---to classify things in order to make sense of them .
    \item We use a \textbf{stock of knowledge} to make sense of the world .
\end{enumerate}

\begin{ExampleBox}
\begin{itemize}
    \item \textbf{Agreed value of money:} We agree that a currency note is a fifty or a hundred; otherwise the note has no value at all . It is only the agreement we have bestowed on objects and individuals that makes the world seem to make sense .
    \item \textbf{Family as agreement:} You agree that she is your mother; you agree that he is your father; you agree that this is your house---everything goes by agreement .
    \item \textbf{Reading an event:} When a death is found in certain circumstances, the first reaction is to call it suicide---because all of us have agreed to read it that way . It could equally have been murder, arranged to give the impression of suicide . There is no objective reality in the reading itself .
\end{itemize}
\end{ExampleBox}

\begin{QuickRevision}
\begin{itemize}
    \item Four propositions: objectivity does not exist in the social world; social order is an illusion; we typify phenomena; we use a stock of knowledge .
    \item Order exists only by shared agreement---currency, family, classroom, the meaning of red .
    \item Even the reading of an event (suicide or murder) is an agreed interpretation, not an objective fact .
\end{itemize}
\end{QuickRevision}

\subsection{Objectivity as Illusion: Reading the Crime Data}

\begin{DataFact}
According to the National Crime Records Bureau (NCRB) report of 2020 or 2021, as stated in class, Delhi and Kerala had the highest crime rate . You read this on the front page and take it as objective reality . It is not . The same figure carries entirely different meanings depending on who reads it .
\end{DataFact}

\vspace{6pt}
\noindent
\begin{tabularx}{\textwidth}{p{3.8cm} Y}
\toprule
\rowcolor{AccentDark}
\thead{Reader} & \thead{Interpretation of ``Delhi has highest crime rate''} \\
\midrule
\textbf{The ordinary reader} & Delhi and Kerala are not safe places to live in . \\
\textbf{The state government} & Shows the accuracy, precision and honesty of recording; administration works efficiently and does not conceal facts . \\
\textbf{A political opponent} & Delhi is not a safe place---the figure is used against the administration . \\
\textbf{Under-reporting reading} & Other states are heavily under-reporting crime; Delhi and Kerala register accurately . This is political data, not factual data . \\
\textbf{Feminist reading} & The prior question is: how do you define crime?  \\
\textbf{Marxist reading} & Reflects bourgeois mentality: targets migrant workers; industrialists who exploit and deny minimum wages are never counted as criminals . \\
\bottomrule
\end{tabularx}
\vspace{6pt}

\paragraph{The feminist question: What counts as crime?}
\begin{itemize}
    \item Crime is itself a social phenomenon . In India the concept of marital rape does not exist in law; in Western countries, where individualism is stronger, marital rape is recognised as a crime .
    \item The difference rests on how marriage is conceived: in Western countries marriage is regarded as a contract, and anything which violates that contract amounts to crime . In India marriage is regarded as a sacred institution---a bond of seven generations, made in heaven .
    \item So when such an act occurs within a marriage in India, the family and even the officials treat it as no crime at all---and because it is not considered a crime, it is never reported, and therefore never enters the crime data of Delhi or Kerala .
\end{itemize}

\begin{CriticalPoint}
\begin{itemize}
    \item \textbf{Uma Chakravarti is cited:} When a woman is in the custody of police officials and they molest or sexually abuse her, it is regarded as a crime . When the same act occurs in the household, it is not regarded as crime . The act is the same . The victim is the same . Only the players have changed---in-laws in place of police officials . How, then, do you define crime? 
    \item \textbf{Law is not objective reality:} Law evolves . Marital rape may be recognised as an offence twenty years from now . There was a point of time when homosexuality was an offence; today it is not . Phenomenology therefore says: objectivity does not exist . We somehow make things objective by forcefully accepting that they exist; otherwise nothing exists as such .
\end{itemize}
\end{CriticalPoint}

\begin{QuickRevision}
\begin{itemize}
    \item The same crime figure yields five different readings---ordinary, governmental, oppositional, feminist, Marxist .
    \item Feminist question: how is crime defined? Marital rape unrecognised in India (sacred bond) vs. West (contract)---never enters data .
    \item Uma Chakravarti: identical act, identical victim; only players change---custody vs. household .
    \item Law itself evolves and is not objective .
\end{itemize}
\end{QuickRevision}

\subsection{Ethnomethodology --- Harold Garfinkel}
Ethnomethodology was formulated by the American sociologist \textbf{Harold Garfinkel} .

\begin{KeyConcept}
\textbf{Ethnomethodology}---the study of the methods used by people to make sense of their everyday lives .
\begin{itemize}
    \item \textit{Ethno} means people; \textit{methodology} means methods (simplified for general understanding) .
\end{itemize}
\end{KeyConcept}

\begin{ExampleBox}
\textbf{Garfinkel's classroom project (The Lodger Experiment):}
\begin{itemize}
    \item Garfinkel told his students to go home and act as an outsider / boarder---not as a son or daughter to their own parents .
    \item A student does so . His mother offers dinner; he eats, and upon finishing asks: \textit{``Nice; how much should I pay? What is the bill?''} 
    \item The mother is shocked and asks whether he has gone mad; she says she will not allow him to go back to the city . She is desperate and anxious to restore the order---the norm that you do not ask your mother for a bill .
    \item The next day students reported: one had been slapped; another's mother was very upset and cried . Garfinkel's response: \textit{``Very good---these are the symbols showing that the mother is desperate to restore what she thinks is the social order of the family.''} 
    \item \textbf{Why Garfinkel did this:} To reveal that what you think of as an objective social order does not really exist---and that even if you tamper with it, the order will be aggressively restored by individuals themselves .
\end{itemize}
\end{ExampleBox}

\begin{ExampleBox}
\textbf{The teacher's phone call breach:} A student called asking for an explanation of ethnomethodology . On the call he was asked: \textit{``Which pizza would you like to have---medium or large?''}  He apologised for a wrong number and hung up . He called again; this time he was asked: \textit{``Where have you been, man? Let us go out for drinks.''}  The student had an internal order of how a call with a teacher proceeds (greeting, decorum, question) . That order was disturbed, and he tried to restore it by concluding he dialled a wrong number . This order has no objective reality; we construct it daily .
\end{ExampleBox}

\begin{KeyConcept}
\begin{itemize}
    \item Every day we construct the family . The family does not exist as an independent entity . Every time you ask your mother how her health is or speak respectfully to your father, you construct family .
    \item \textbf{Dictated core:} \textit{Ethnomethodology agrees with phenomenology that social order is an illusion. However, ethnomethodology performs experiments to reveal the chaos underlying the social order. It is, in that sense, more experimental phenomenology.} 
\end{itemize}
\end{KeyConcept}

\begin{QuickRevision}
\begin{itemize}
    \item Harold Garfinkel---ethnomethodology = study of methods people use to make sense of everyday lives .
    \item The lodger/bill experiment: the order is broken, and the other person becomes anxious to restore it .
    \item Agrees with phenomenology that order is illusion, but performs experiments to reveal the chaos beneath it .
    \item Family does not exist as a thing; it is constructed daily through performances .
\end{itemize}
\end{QuickRevision}

\subsection{The Experiments: Breaching, Conversation Analysis, Documentary Method}
Ethnomethodologists use specific methods to reveal how order is maintained :
\begin{enumerate}
    \item \textbf{Breaching experiment:} Deliberately disrupting standard social expectations to observe how actors react to restore order .
    \begin{itemize}
        \item \textit{Tic-tac-toe:} An opponent places a mark outside the grid lines, breaching the order; you restore order by calling him mad or illiterate .
        \item \textit{The bookshop:} Asking for ten successive book titles, having the shopkeeper bring each, and then walking away without buying; the shopkeeper restores order through angry protest .
    \end{itemize}
    \item \textbf{Conversation analysis:} Examining the micro-structure of spoken talk and sequential turn-taking that make conversations appear normal (e.g., greetings, pauses, tone) .
    \begin{itemize}
        \item \textit{Unlaughed joke:} Cracking a joke where the listener refuses to laugh; the breach reveals that laughing is an expected social convention, not an objective necessity .
    \end{itemize}
    \item \textbf{Documentary method:} Members of society select specific aspects of a situation, interpret them as evidence of an underlying pattern (a ``document''), and then use that pattern to interpret all subsequent actions .
    \begin{itemize}
        \item \textit{The officer making a visitor wait:} An administrative officer sits idle but asks an assistant to make a visitor wait for 25 minutes . Why? Because the mental ``document'' of bureaucracy dictates that an officer must look busy and inaccessible . We reproduce these performances to sustain the appearance of an officer's authority .
    \end{itemize}
\end{enumerate}

\begin{QuickRevision}
\begin{itemize}
    \item Two experimental tools: breaching experiments and conversation analysis .
    \item Documentary method: actors carry internal scripts/templates from socialisation and reproduce them .
    \item Illustrations: the bill, tic-tac-toe, the bookshop, the unlaughed joke, the officer making a visitor wait .
\end{itemize}
\end{QuickRevision}

\subsection{Symbolic Interactionism}
The lines to be written for symbolic interactionism :
\begin{enumerate}
    \item The world consists of \textbf{symbols}, and individuals make sense of the world by interpreting the symbols . Everything is symbol; there is no reality apart from it .
    \item Social order is the result of individuals interpreting the symbols in a \textbf{uniform fashion} .
    \item A symbol in itself does not have any inherent meaning; rather, we as \textbf{active players attach meaning} to those symbols .
\end{enumerate}

\begin{ExampleBox}
\begin{itemize}
    \item \textbf{Classroom objects:} The laptop, pen, notebook, books, table, chair, T-shirt, watch, hairstyle, and hand gestures are all symbols that we collectively decode .
    \item \textbf{Gated bungalow:} A huge gate is a symbol signifying: \textit{stay away, elites reside here} .
    \item \textbf{Advertisements:} An instant noodle or chocolate ad appears while studying, symbolising comfort/craving, triggering consumption .
    \item \textbf{Study material:} Buying a specific book because it symbolises success in clearing the examination .
\end{itemize}
\end{ExampleBox}

\begin{CriticalPoint}
\begin{itemize}
    \item \textbf{Meanings of symbols are contextual:} A thumbs-up gesture that signifies ``all the best'' in India may be read as an obscene insult in parts of Russia or the Middle East .
    \item \textbf{George Herbert Mead:} Formulated Symbolic Interactionism .
    \item \textbf{Tradition grouping:} Symbolic interactionism and ethnomethodology are phenomenological schools . All three reject objective social structures, demonstrating that order is constructed daily via symbols and mutual expectations .
\end{itemize}
\end{CriticalPoint}

\begin{QuickRevision}
\begin{itemize}
    \item World consists of symbols; individuals make sense by interpreting them .
    \item Social order = symbols interpreted in a uniform fashion .
    \item Symbols have no inherent meaning---active players attach it; meaning is contextual .
    \item Given by George Herbert Mead; SI and ethnomethodology are phenomenological schools .
\end{itemize}
\end{QuickRevision}

\subsection{One Phenomenon, Five Perspectives: Two People Decide to Marry}
To demonstrate the analytical distinctions between perspectives, the same social phenomenon---A falls in love with B and decides to marry---is examined across five frameworks :

\paragraph{1. Auguste Comte --- Positivism}
\begin{itemize}
    \item It was natural, and a part of the evolutionary process .
    \item Marriage forms families, and families are the units responsible for the maintenance of order .
    \item The two individuals are simply passive robots forming a family, fulfilling the need of the larger society to produce tomorrow's labour force and maintain social stability .
\end{itemize}

\paragraph{2. Max Weber --- Social Action}
\begin{itemize}
    \item Rejects overarching laws; focuses on individual \textbf{meanings and motives} .
    \item The man marries this woman because she is the daughter of a high-ranking official, providing social capital, business networks, and political entry . The woman seeks security and status preservation .
    \item Love is a subjective meaning constructed by active actors pursuing utilitarian and status goals .
\end{itemize}

\paragraph{3. Alfred Schutz --- Phenomenology}
\begin{itemize}
    \item Love, marriage, and family are complete illusions with no objective reality .
    \item The sacredness of marriage is manufactured through external rituals: exchanging rings, booking wedding halls, printing invitation cards, and lavish catering expenditure .
    \item It is the shared subjective belief that one is married that creates the appearance of reality .
\end{itemize}

\paragraph{4. Harold Garfinkel --- Ethnomethodology}
\begin{itemize}
    \item Marriage is an ongoing, fragile interactional construction .
    \item \textit{Breaching illustration:} A student stages a complete one-month marriage ceremony as a class project, and at the end of the reception informs the bride it was merely an experiment for 300 marks . The resulting tears, shock, and parental interventions reveal the desperate social machinery deployed to restore the illusion of order .
\end{itemize}

\paragraph{5. George Herbert Mead --- Symbolic Interactionism}
\begin{itemize}
    \item Marriage is sustained through an interconnected web of decoded symbols .
    \item The wedding invitation card, the DJ music, the decorated pandal, the stage lighting, the couple's welcome photograph, the banquet food, the \textit{sindoor}, and the \textit{mangalsutra} are symbols .
    \item When hundreds of guests agree on what these symbols represent, the orderly social phenomenon of a ``wedding'' emerges .
\end{itemize}

\begin{QuickRevision}
\begin{itemize}
    \item \textbf{Comte (Positivism):} Evolutionary necessity, family reproduces labour for social maintenance .
    \item \textbf{Weber (Social Action):} Subjective motives (status, political connections, utilitarian security) .
    \item \textbf{Schutz (Phenomenology):} Manufactured illusion via rituals, ring exchanges, and shared belief .
    \item \textbf{Garfinkel (Ethnomethodology):} Fragile social order revealed by breaching matrimonial expectations .
    \item \textbf{Mead (Symbolic Interactionism):} Uniform decoding of symbols (\textit{sindoor}, \textit{mangalsutra}, pandal, cards) .
\end{itemize}
\end{QuickRevision}

\subsection{Dramaturgy --- Erving Goffman}

\begin{KeyConcept}
\textbf{Dramaturgy (Erving Goffman):} As Shakespeare wrote, \textit{``All the world's a stage, and all the men and women merely players.''} 
\begin{itemize}
    \item Social interaction is analyzed as a theatrical performance where actors follow scripts, adhere to role expectations, and adjust behaviour to specific settings .
    \item \textbf{The Proposal as Drama:} A young man seeks advice from friends on proposing . He is given a pre-written script (e.g., \textit{``The moment you entered my life, everything changed''}) . He takes her to dinner and delivers his lines . She acts her part (blushing, replying \textit{``How cute''}) . Both fulfill culturally scripted role expectations .
    \item Individuals are active players who switch roles across contexts (nervous lover $\rightarrow$ disciplined student $\rightarrow$ authoritative district magistrate) . We are not robots; we hold our own remote control .
\end{itemize}
\end{KeyConcept}

\begin{CriticalPoint}
\begin{itemize}
    \item These perspectives are analytical frameworks for the exam, not personal philosophies to turn you into a cynic or nihilist .
    \item Avoid directly copying classroom examples; cultivate the skill of observing your own everyday surroundings through these lenses .
\end{itemize}
\end{CriticalPoint}

\begin{QuickRevision}
\begin{itemize}
    \item Dramaturgy---Erving Goffman; theatrical model of social interaction .
    \item Scripted performances and contextual role switching prove individuals are active agents, not passive robots .
\end{itemize}
\end{QuickRevision}

\subsection{Conclusion: Grounds on Which Non-Positivism Critiques Positivism}

\begin{KeyConcept}
\textbf{Dictated conclusion:} \textit{In conclusion, non-positivism critiques positivism on the grounds that individuals are not passive actors or governed by laws. Rather, they are involved in:} 
\begin{enumerate}
    \item Ascribing \textbf{meanings and motives} to their actions---(\textbf{Max Weber}) ;
    \item Interpreting the \textbf{meaning of symbols} to make sense of their everyday lives---(\textbf{George Herbert Mead / Symbolic Interactionism}) ;
    \item Performing their roles by doing \textbf{impression management}---(\textbf{Erving Goffman}) ;
    \item Creating an \textbf{orderly pattern} which they feel is required to make social order possible in that context .
\end{enumerate}
\end{KeyConcept}

\paragraph{Impression Management and Contextual Order}
\begin{itemize}
    \item \textbf{Front stage vs. Backstage:} Asking friends for proposal tips was backstage preparation; delivering the proposal with confidence was front stage performance .
    \item Because an individual performs different roles before different audiences, a hundred people can hold a hundred distinct opinions about the same individual .
    \item \textit{The airplane wedding breach:} During the lockdown, a couple married inside an airplane, violating the conventional ``document'' of weddings, creating media sensation until public discourse rationalised it .
    \item Positivism describes only structural functions and passive recipients; non-positivism centers human agency and meaning construction .
\end{itemize}

\begin{PYQLink}
These perspectives are essential for UPSC questions on \textit{Sociology vs. Common Sense}, \textit{Sociology vs. Psychology}, and the core methodology debates in Unit 2 and Unit 4 .
\end{PYQLink}

\begin{QuickRevision}
\begin{itemize}
    \item Non-positivism critiques positivism via: Weber (motives), Mead (symbols), Goffman (impression management), and phenomenologists (contextual pattern creation) .
    \item Positivism = macro-structures and passive robots; Non-positivism = micro-agency and active meaning-makers .
\end{itemize}
\end{QuickRevision}

\subsection{Doubts Discussed in Session 4}

\paragraph{1. How do Schutz and Weber relate in the marriage example?}
\begin{itemize}
    \item Weber introduced social action; Schutz developed phenomenology out of social action . Schutz regarded Weber's work as indispensable for breaking positivism's obsession with structural laws .
    \item \textbf{Weber} focuses on the actor's \textbf{intentions} (meanings and motives---why this specific person marries for status or capital) .
    \item \textbf{Schutz} goes deeper, investigating the \textbf{epistemological construction of meaning} (how marriage is taken for granted as a real institution via shared typifications) .
\end{itemize}

\begin{ExampleBox}
\textbf{Arranged marriage ritual:} A matchmaking meeting involves tea, snacks, polite family queries, and references to district work . Inside the culture, the ritual makes perfect sense; to an outsider unfamiliar with the codes, it appears as an inexplicable violation of personal privacy . The reality of the meeting is entirely socially constructed .
\end{ExampleBox}

\paragraph{2. Further breaching examples discussed}
\begin{ExampleBox}
\begin{itemize}
    \item \textbf{National anthem in a fast-food restaurant:} Playing the national anthem while people are eating burgers breaches restaurant decorum, causing acute confusion .
    \item \textbf{National anthem in movie theatres:} Audiences stand automatically, not necessarily out of spontaneous devotion, but because standing is enforced by social expectations and fear of being labeled deviant .
    \item \textbf{Reversing traffic lights:} If everyone stopped on green and moved on red, instant panic and accidents would follow until actors aggressively tried to re-establish the normative rule .
\end{itemize}
\end{ExampleBox}

\begin{QuickRevision}
\begin{itemize}
    \item Weber = intentions/motives; Schutz = construction of shared meaning .
    \item Breaching examples (anthem in restaurant/cinema, reversing traffic lights) confirm that social order is maintained by active convention, not natural law .
\end{itemize}
\end{QuickRevision}

\sectiondivider

\section{Applying the Perspectives, Positivism in Depth and Its Critique}

\subsection{Applying the Perspectives I --- Is the Nation Real or an Illusion?}
For phenomenology and ethnomethodology, the concept of the nation is a socially constructed illusion . The nation does not exist as a physical object; it exists because we continuously perform and accept its symbols .

\begin{ExampleBox}
\textbf{Constructing Indian national identity:}
\begin{itemize}
    \item In primary school, you attend morning assembly at 8:30 a.m., singing the national anthem listing Punjab, Sindh, Gujarat, Maratha .
    \item In history class, you are taught about Gandhi, Nehru, Ambedkar, and the freedom struggle .
    \item Geography maps demarcate territorial boundaries . By age 25, the nation feels like an absolute, concrete reality . The same construction applies to the USA, Israel, or Palestine .
\end{itemize}
\end{ExampleBox}

\begin{KeyConcept}
\textbf{Banal Nationalism (Michael Billig):} The everyday, routine habits, language, national anthems, flag hoisting, and representations that continuously reproduce national consciousness .
\end{KeyConcept}

\begin{CriticalPoint}
\textbf{George Orwell's \textit{Animal Farm}:} \textit{``All animals are equal, but some animals are more equal than others.''}  In formal theory, constitutional democracy promises equal opportunity; in practice, privileged networks dominate . Aspirants competing for civil services are striving to become ``more equal'' within the formal system .
\end{CriticalPoint}

\begin{QuickRevision}
\begin{itemize}
    \item Nation is a socially constructed illusion sustained by repetition .
    \item Daily rituals (anthems, maps, history lessons) generate \textbf{banal nationalism} .
\end{itemize}
\end{QuickRevision}

\subsection{Applying the Perspectives II --- Sex, Gender and Performance}

\begin{KeyConcept}
\begin{itemize}
    \item \textbf{Sex is biological and real; Gender is socially constructed and performative.} 
    \item These theories do not deny your physical body or biological organs; they deny that the cultural meanings attached to them are natural or fixed .
\end{itemize}
\end{KeyConcept}

\begin{ExampleBox}
\textbf{Performing Masculinity and Femininity:}
\begin{itemize}
    \item \textit{The crowded bakery:} Shouting aggressively to get bread displays masculine dominance .
    \item \textit{Festival greetings:} Speaking gently and warmly displays emotional traits culturally coded as feminine .
    \item \textit{A woman riding a heavy motorcycle:} Displaying tough, assertive behaviour disrupts traditional gender scripts .
\end{itemize}
\end{ExampleBox}

\begin{KeyConcept}
\textbf{Gender as a Fluid Performance (Judith Butler):}
Gender is not an innate essence . It is a continuous performative spectrum along which individuals toggle . \textbf{It is not your gender that determines your performance; your performance decides your gender.} 
\end{KeyConcept}

\begin{ExampleBox}
Indira Gandhi was described as \textit{``the only man in her cabinet''} because her assertive decision-making aligned with traits culturally reserved for men .
\end{ExampleBox}

\begin{QuickRevision}
\begin{itemize}
    \item Sex = biological fact; Gender = performative spectrum .
    \item Performance determines gender, not biology .
    \item Nation, gender, family, and religion are continuously sustained through repetitive performance .
\end{itemize}
\end{QuickRevision}

\subsection{Positivism: The Philosophy and Its Laws}

\begin{KeyConcept}
\textbf{Positivism} is the epistemological philosophy which asserts that society, individuals, groups, and human thought are governed by invariant \textbf{social laws} analogous to natural laws .
\begin{itemize}
    \item For pure positivism, individual consciousness is secondary; individuals are components of a larger organism governed by structural laws they cannot alter .
    \item Just as gravity operates invisibly in physics, social laws govern historical progress and social integration .
\end{itemize}
\end{KeyConcept}

\subsection{The Law of Three Stages}
Auguste Comte formulated the \textbf{Law of Three Stages}, arguing that human thought and social organisation pass through three sequential epochs :

\vspace{6pt}
\noindent
\begin{tabularx}{\textwidth}{p{3.2cm} p{3.2cm} Y}
\toprule
\rowcolor{AccentDark}
\thead{Stage} & \thead{Period} & \thead{Mode of Thinking \& Social Structure} \\
\midrule
\textbf{Theological Stage} & Up to c.~1300 AD & \textbf{Religious explanations:} Natural and social events attributed to gods, spirits, or divine will; feudal hierarchy and absolute monarchy legitimated by divine order . \\
\textbf{Metaphysical Stage} & c.~1300--1800 AD & \textbf{Abstract philosophical concepts:} Transition era; nature, reason, and abstract rights replace personified gods, but religion persists alongside early science . \\
\textbf{Positive (Scientific) Stage} & c.~1800 onward & \textbf{Empirical and scientific laws:} Explanations based on observation, experimentation, and invariant laws; emergence of industrial society and sociology . \\
\bottomrule
\end{tabularx}
\vspace{6pt}

\begin{QuickRevision}
\begin{itemize}
    \item Law of Three Stages: Theological ($-\text{1300}$) $\rightarrow$ Metaphysical ($1300\text{--}1800$) $\rightarrow$ Positive ($1800-$) .
    \item The law determines both intellectual modes of thought and social structures .
\end{itemize}
\end{QuickRevision}

\subsection{Social Statics and Social Dynamics}
Comte divided sociology into two foundational branches :
\begin{enumerate}
    \item \textbf{Social Statics (Law of Coexistence):} The study of the conditions of order, stability, and harmony among social institutions (religion, family, state, property) .
    \item \textbf{Social Dynamics (Law of Succession):} The study of the laws of social progress, historical evolution, and systemic transformation .
\end{enumerate}

\paragraph{The Historical Dialectic of Statics and Dynamics}
\begin{itemize}
    \item \textit{Up to 1300 AD:} Social statics dominated (stable theological feudal order) .
    \item \textit{1300 to 1800 AD:} Intellectual crises (Renaissance, scientific revolution, Enlightenment) created institutional friction, initiating social dynamics .
    \item \textit{Revolutions (1750--1790s):} Industrial and French Revolutions destroyed feudal statics, producing temporary anarchy .
    \item \textit{1800 onward:} Emergence of the positive stage restoring social statics on a scientific foundation .
\end{itemize}

\begin{CriticalPoint}
\begin{itemize}
    \item Comte's framework is an analytical evolutionary model, not an assertion that traditional inequalities (caste, patriarchy) have vanished empirically .
    \item Modern societies focus on systemic reforms within constitutional frameworks rather than total violent revolutions .
    \item Social change is an evolutionary wave governed by structural laws, not the creation of isolated individuals .
\end{itemize}
\end{CriticalPoint}

\begin{QuickRevision}
\begin{itemize}
    \item Social Statics = order/coexistence; Social Dynamics = progress/evolution .
    \item Historical cycle: Feudal statics $\rightarrow$ Enlightenment dynamics $\rightarrow$ Revolutionary crisis $\rightarrow$ Positive scientific order .
\end{itemize}
\end{QuickRevision}

\subsection{Prediction: The Power Comte Claimed for Sociology}

\begin{KeyConcept}
\textbf{Comte's Maxim:} \textit{``From science comes prevision; from prevision comes action.''} 
\begin{itemize}
    \item Scientific discovery of social laws enables sociology to \textbf{predict} future social developments and engineer institutional reforms to eliminate social pathologies (poverty, crime, war, inequality) .
    \item \textit{Meteorological analogy:} Just as a meteorologist predicts a cyclone to save lives through planned evacuations, the sociologist predicts social crises to restore equilibrium .
\end{itemize}
\end{KeyConcept}

\begin{QuickRevision}
\begin{itemize}
    \item Positivism claims predictive capacity for sociology to guide social engineering and resolve crises .
\end{itemize}
\end{QuickRevision}

\subsection{Applying Scientific Method to Society --- ``We Are Apples''}

\begin{CriticalPoint}
\begin{itemize}
    \item Positivism assumes the methodology of natural sciences applies directly to human society .
    \item \textbf{The fundamental flaw:} An apple falling from a tree does not have consciousness, reflexivity, or subjective intent; it is governed purely by external gravity .
    \item Positivism treats human beings as ``apples''---passive objects governed by external social forces, ignoring consciousness and agency .
\end{itemize}
\end{CriticalPoint}

\begin{KeyConcept}
\textbf{Immanuel Kant on Enlightenment:} \textit{``Sapere Aude''}---Have the courage to use your own reason and intelligence rather than blindly following traditional authorities or external dogmas .
\end{KeyConcept}

\begin{QuickRevision}
\begin{itemize}
    \item Positivism treats human beings as ``apples'' governed by external laws .
    \item Ignores human consciousness, meaning, and reflexivity .
    \item Kant: Enlightenment is thinking autonomously with your own reason .
\end{itemize}
\end{QuickRevision}

\subsection{Positivism, Non-Positivism and Anti-Positivism}

\begin{KeyConcept}
\textbf{Definitional Distinctions for Answer Writing:}
\begin{itemize}
    \item \textbf{Non-Positivist Methodology (Singular):} Refers specifically to \textbf{Max Weber's Social Action Theory} .
    \item \textbf{Non-Positivist Methodologies (Plural) / Anti-Positivism:} Refers to \textbf{Phenomenology, Ethnomethodology, Symbolic Interactionism, and Dramaturgy} .
\end{itemize}
\end{KeyConcept}

\vspace{6pt}
\noindent
\begin{tabularx}{\textwidth}{p{4.2cm} p{4.2cm} Y}
\toprule
\rowcolor{AccentDark}
\thead{Perspective} & \thead{Do individuals exist?} & \thead{Do structures exist?} \\
\midrule
\textbf{Positivism} & No (passive recipients) & Yes (real and law-governed)  \\
\textbf{Non-Positivism} (Max Weber) & Yes (actors have meanings/motives) & Yes (structures arise as consequences of social actions)  \\
\textbf{Anti-Positivism} (Phenomenology, Ethnomethodology, SI, Dramaturgy) & Yes (active, manipulative agents) & No (structures are socially constructed illusions)  \\
\bottomrule
\end{tabularx}
\vspace{6pt}

\begin{CriticalPoint}
Statements like ``nation is illusion'' or ``gender is illusion'' are specific anti-positivist perspectives, not absolute ontological truths . Sociology is multi-paradigmatic .
\end{CriticalPoint}

\begin{QuickRevision}
\begin{itemize}
    \item Positivism = structures only; Non-positivism (Weber) = individuals + structures; Anti-positivism = individuals only (structures are illusions) .
\end{itemize}
\end{QuickRevision}

\subsection{Critique of Positivism and the Opening of Critical Theory}
Non-positivism is not the sole critique of positivism . A complete evaluation requires introducing \textbf{Critical Theory} (the Frankfurt School) :
\begin{itemize}
    \item \textbf{First-Generation Critical Theorists:} \textbf{Theodor Adorno} and \textbf{Max Horkheimer} .
    \item \textbf{Core Critique:} Positivism is not value-neutral science; it functions as an ideology of domination that justifies the status quo, reduces humans to instruments, and legitimates technocratic control .
    \item \textit{Note on generations:} \textbf{Jürgen Habermas} belongs to the second generation of critical theorists (not the first) .
\end{itemize}

\begin{QuickRevision}
\begin{itemize}
    \item Critical Theory (Adorno and Horkheimer---1st generation) critiques positivism as an ideology of domination .
    \item Habermas = 2nd generation critical theorist .
\end{itemize}
\end{QuickRevision}

\sectiondivider

\section{Answer-Writing Toolkit}

\subsection{How to Attack the Question}
\begin{enumerate}
    \item \textbf{Encircle Keywords First:} In \textit{``Is sociology common sense?''}, identify \textit{sociology}, \textit{common sense}, and \textit{is} .
    \item \textbf{Define Keywords Succinctly:} Provide brief, crisp definitions suited to the marks .
    \item \textbf{Establish the Analytical Relationship:} Connect the defined terms before declaring an evaluative verdict .
    \item \textbf{Maintain Sequential Flow:} Ensure every sentence logically proceeds from the preceding one .
    \item \textbf{Use ``However'' for Dialectical Pivot:} Force a balanced transition to counter-perspectives using ``However'' .
    \item \textbf{Deliver a Mature, Balanced Conclusion:} Avoid flat ``yes'' or ``no'' verdicts; characterize shifts as the intellectual maturity of the discipline .
\end{enumerate}

\subsection{Presentation, Length and Language}
\begin{itemize}
    \item \textbf{Word Limit Discipline:} 10 marks $\approx 150$ words; 20 marks $\approx 250$ words .
    \item \textbf{Underline Keywords:} Use your standard blue pen only (no red, pink, or black pens) .
    \item \textbf{Form Over Content:} Structural clarity, logical scaffolding, and legible presentation yield high marks .
    \item \textbf{Formula for Success:} $\text{Limited Information} + \text{Unlimited Revision} = \text{Clarity of Thought}$ .
    \item \textbf{Administrative Neutrality:} Avoid naming living political figures; attribute institutional statements to bodies (e.g., \textit{``The Indian Medical Association classified...''}) .
\end{itemize}

\subsection{Exam Framing, Sources and Statements to Remember}
\begin{itemize}
    \item \textit{Sociology vs. Philosophy:} Never asked by UPSC (Philosophy is the mother of all sciences) .
    \item \textit{Sociology vs. Political Science:} Now evaluated through the lens of interdisciplinarity .
    \item \textit{Sociology vs. Anthropology:} Refer to NCERT Class 11 \textit{Introducing Sociology}, Chapter 1 .
    \item \textbf{Key Exam Statement:} \textit{``Science is a social process''} (Thomas Kuhn) .
    \item \textbf{Terminology Rule:} Non-positivist methodology (singular) = Max Weber; Non-positivist methodologies (plural) = Anti-positivism (Phenomenology, Ethnomethodology, SI, Dramaturgy) .
    \item Use the philosophy and history of science (Popper, Kuhn, Feyerabend) to construct high-scoring philosophical essays in UPSC Mains .
\end{itemize}